%%%%%%%%%%%%%%%%%%%%%%%%%%%%%%%%%%%%%%%%%%%%%%%%%%%%%%%%%%%%%%%%%%%%%%
% Lớp 11 | Dãy số. Cấp số cộng. Cấp số nhân | Dãy số
% Dạng 1D2?1-4: Dãy số tăng, dãy số giảm
% Số câu: 107  (sinh tự động bởi tach_cau.py; từ nay thầy sửa trực tiếp file này)
%%%%%%%%%%%%%%%%%%%%%%%%%%%%%%%%%%%%%%%%%%%%%%%%%%%%%%%%%%%%%%%%%%%%%%

%%%%%%%%%%%%%%%% TRẮC NGHIỆM 4 PHƯƠNG ÁN %%%%%%%%%%%%%%%%

% nguồn: Ok_Toan11_TracNghiem_Lan2
\begin{ex}%[1D2N1-4]%[ma:75f00a73]
	Cho các dãy số sau, dãy số  nào không là dãy số tăng?
	 \choice
		{\True  $1; 1; 1; 1;\cdots$}
		 {$1; 3; 5; 7;\cdots$}
		{$2; 4;  6; 8;\cdots$}
		{$\dfrac{1}{2}; 1; \dfrac{3}{2}; 2;\cdots$} 
	 \loigiai{
		Xét đáp án $1; 1; 1; 1;\cdots$ ta có dãy $1; 1; 1; 1;\cdots$ là dãy hằng nên không  tăng không giảm.
	} 
\end{ex}

% nguồn: Ok_Toan11_TracNghiem_Lan2
\begin{ex}%[1D2N1-4]%[ma:adac72e6]
	Dãy số nào dưới đây là dãy  số tăng?
	 \choice
		{$u_n=-2n+1$} 
		{\True $v_n=3n-10$} 
		{$a_n=\dfrac{1}{n}$} 
		 {$b_n=(-1)^n$}
	\loigiai{ 
		Xét $u_n=-2n+1$. Ta có $u_1=-1> u_2=-3$, suy ra dãy không tăng.\\ 
		Xét $a_n=\dfrac{1}{n}$. Ta có $a_1=1> a_2=\dfrac{1}{2}$,  suy ra dãy không tăng.\\
		Xét $b_n=(-1)^n$. Ta có $b_2=1> b_3=-1$, suy ra dãy  không tăng.\\
		Xét $v_n=3n-10$. Ta có $v_{n+1}-v_n=  3(n+1)-10-(3n-10)=3$, suy ra dãy số tăng.\\
		 Vậy chọn dãy số $v_n=3n-10$.
	} 
\end{ex}

% nguồn: Ok_Toan11_TracNghiem_Lan2
\begin{ex}%[1D2N1-4]%[ma:9e428549]
	Trong các dãy số sau, dãy số nào là dãy số  giảm.
	\choice 
		 {\True $\dfrac{1}{2};\dfrac{1}{4};\dfrac{1}{6};\dfrac{1}{8};\cdots;\dfrac{1}{2n};\cdots$}
		 {$6;9;12;15;\cdots;3n+3;\cdots$}
		{$2;4;8;16;\cdots;2^n;\cdots$} 
		{$-5;-5;-5;-5;-5;\cdots;-5;\cdots$} 
	 \loigiai{
		Ta có dãy  số $\dfrac{1}{2};\dfrac{1}{4};\dfrac{1}{6};\dfrac{1}{8};\cdots;\dfrac{1}{2n};\cdots$ giảm.\\
		 Dãy số $6;9;12;15;\cdots;3n+3;\cdots$ tăng.\\
		Dãy số $2;4;8;16;\cdots;2^n;\cdots$  tăng.\\
		Dãy số $-5;-5;-5;-5;-5;\cdots;-5;\cdots$ không đổi. 
	 }
\end{ex}

% nguồn: Ok_Toan11_TracNghiem_Lan2
\begin{ex}%[1D2N1-4]%[ma:a3303c87]
	Cho các dãy số sau.  Dãy số nào \textbf{không} là dãy số tăng?
	\choice 
		{\True  $1; 1; 1; 1; \ldots$}
		{$1; 3; 5;  7; \ldots$}
		{$2; 4; 6; 8; \ldots$} 
		{$\dfrac{1}{2}; 1;  \dfrac{3}{2}; 2; \ldots$}
	\loigiai{ 
		Ta có  dãy $1; 1; 1; 1; \ldots$ là dãy hằng nên không tăng không giảm.
	 }
\end{ex}

% nguồn: Ok_Toan11_TracNghiem_Lan2
\begin{ex}%[1D2N1-4]%[ma:10d95637]
	Cho dãy số $(u_n)$ biết $u_n=\left(-1\right)^n \left(n^2+1\right)$. Mệnh đề nào  sau đây đúng?
	\choice 
		{Dãy số tăng} 
		{Dãy  số giảm}
		{\True Dãy số  không tăng, không giảm}
		 {Dãy số là dãy hữu hạn}
	 \loigiai{
		Dãy không tăng, không giảm  vì các số hạng đan dấu.
	} 
\end{ex}

% nguồn: Ok_Toan11_TracNghiem_Lan2
\begin{ex}%[1D2N1-4]%[ma:d9b189c1]
	Cho dãy số $(u_n)$, biết $u_n=n+1$. Khẳng định nào sau đây  là đúng?
	\choice 
		{\True Dãy số $(u_n)$ là dãy số  tăng}
		{Dãy số $(u_n)$ là dãy  số giảm}
		{Dãy số $(u_n)$ không là dãy số tăng, cũng không là  dãy số giảm}
		{Dãy số $(u_n)$ là dãy số  bị chặn}
	\loigiai{ 
		Ta có $u_{n+1}-u_n=(n+2)-(n+1)=1> 0$, tức là $u_{n+1} > u_n$ với $\forall  n\in \mathbb{N}^{*}$.\\
		Vậy dãy số $(u_n)$ là dãy số  tăng.
	} 
\end{ex}

% nguồn: Ok_Toan11_TracNghiem_Lan2
\begin{ex}%[1D2N1-4]%[ma:aeb64f41]
	Cho các dãy số sau, dãy số nào là dãy  số tăng?
	\choice 
		{$1; 1;  1; 1; 1; 1;\cdots$}
		{$1; -\dfrac{1}{2}; \dfrac{1}{4}; -\dfrac{1}{8};  \dfrac{1}{16};\cdots$}
		{\True $1; 3; 5; 7;  9;\cdots$}
		{$1;  \dfrac{1}{2}; \dfrac{1}{4}; \dfrac{1}{8}; \dfrac{1}{16};\cdots$}
	 \loigiai{
		Xét đáp án $1; 1; 1; 1; 1; 1;\cdots$ đây là dãy hằng nên dãy  không tăng, không giảm.\\
		Xét  đáp án $1; -\dfrac{1}{2}; \dfrac{1}{4}; -\dfrac{1}{8}; \dfrac{1}{16};\cdots \Rightarrow u_1 > u_2 < u_3 \Rightarrow $ loại.\\
		Xét  đáp án $1; 3; 5; 7; 9;\cdots \Rightarrow u_n < u_{n+1}, n\in \mathbb{N}^{*}$.\\
		Xét đáp án $1; \dfrac{1}{2}; \dfrac{1}{4}; \dfrac{1}{8}; \dfrac{1}{16};\cdots \Rightarrow u_1 >  u_2 > u_3 \cdots > u_n > \cdots$, suy ra loại.
	 }
\end{ex}

% nguồn: Ok_TracNghiemLop11_Lan1
\begin{ex}%[1D2N1-4]%[ma:dcfd9246]
	Cho các dãy số sau. Dãy số nào là dãy số tăng? 
	\choice 
		{$1$;  $1$; $1$; $1$; $1$; $1;\ldots$}
		 {$1$; $-\dfrac{1}{2}$; $\dfrac{1}{4}$; $-\dfrac{1}{8}$; $\dfrac{1}{16};\ldots$}
		{\True $1$; $3$; $5$;  $7$; $9;\ldots$}
		{$1$; $\dfrac{1}{2}$; $\dfrac{1}{4}$; $\dfrac{1}{8}$; $\dfrac{1}{16};\ldots$} 
	\loigiai{ 
	Dãy số tăng là dãy số  $1$; $3$; $5$; $7$; $9;\ldots$ vì $u_n<u_{n+1}$, $\forall n \in \mathbb{N}^*$.}
\end{ex}

% nguồn: Ok_TracNghiemLop11_Lan1
\begin{ex}%[1D2N1-4]%[ma:32f6d1f3]
	 Trong các dãy số $(u_{n})$ có số hạng tổng quát dưới đây, dãy số nào là dãy số tăng?
	\choice 
		{$u_{n}=3-n^5$} 
		 {$u_{n}=\dfrac{1}{n}$}
		{\True $u_{n}=2n^3$} 
		 {$u_{n}=\dfrac{5}{n+3}$}
	\loigiai{ 
		Dãy số $u_{n}=3-n^5$ có $u_1=3-1^5=2, u_2=3-2^5=-29 \Rightarrow u_1>u_2$. Đây  không là dãy số tăng.\\
		Dãy số $u_{n}=\dfrac{1}{n}$ có $u_1=\dfrac{1}{1}=1,u_2=\dfrac{1}{2}\Rightarrow u_1>u_2$. Đây không là  dãy số tăng.\\
		Dãy số $u_{n}=2n^3$ có $u_{n+1}=2(n+1)^3=2n^3+6n^2+6n+2, u_n=2n^3 \Rightarrow u_{n+1}>u_n, \forall n\in  \mathbb{N^*}$. Đây là dãy số tăng.\\
		 Dãy sô $u_{n}=\dfrac{5}{n+3}$ có $u_1=\dfrac{5}{1+3}=\dfrac{5}{4}, u_2=\dfrac{5}{2+3}=\dfrac{5}{5}=1 \Rightarrow u_1>u_2$. Đây không là dãy số tăng.
	} 
\end{ex}

% nguồn: Ok_TracNghiemLop11_Lan1
\begin{ex}%[1D2N1-4]%[ma:8fb1eed5]
	Cho dãy số $(u_n)$ có $u_n=\sin \dfrac{\pi}{n}$. Chọn  khẳng định đúng trong các khẳng định sau đây.
	 \choice
		{\True Dãy số $(u_n)$ bị chặn} 
		{Dãy số $(u_n)$ tăng} 
		{Dãy số $(u_n)$ giảm} 
		{Dãy số $(u_n)$ không bị chặn} 
	\loigiai{ 
	Do $-1\le \sin \dfrac{\pi}{n}\le 1$ nên dãy số  $(u_n)$ bị chặn.}
\end{ex}

% nguồn: Ok_TracNghiemLop11_Lan1
\begin{ex}%[1D2N1-4]%[ma:6c3e9020]
	Cho các dãy số sau, dãy số  nào là dãy tăng?
	 \choice{$1$; $-\dfrac{1}{2}$; $\dfrac{1}{3}$; $-\dfrac{1}{4}$}
		{$1$;  $\dfrac{1}{2}$; $\dfrac{1}{4}$; $\dfrac{1}{6}$}
		 {$-1$; $3$; $5$; $3$}
		{\True $2$; $4$;  $6$; $8$}
	\loigiai{Dãy  số tăng là $2$; $4$; $6$; $8$.}
\end{ex}

% nguồn: Ok_Toan11_TracNghiem_Lan2
\begin{ex}%[1D2H1-4]%[ma:4f227422]
	Trong các dãy số  $\left(u_n \right)$ cho bởi số hạng tổng quát $u_n$ sau, dãy số nào tăng?
	 \choice
		{$u_n=\dfrac{1}{3^n}$} 
		{$u_n=\dfrac{1}{2n+1}$} 
		{$u_n=\dfrac{n+1}{3n+2}$} 
		{\True  $u_n=\dfrac{4n-2}{n+3}$}
	 \loigiai{
		Xét $u_n=\dfrac{1}{3^n}$.  Ta có
		$u_{n+1}-u_n=\dfrac{1}{3^{n+1}}-\dfrac{1}{3^n}=\dfrac{1}{3.3^n}-\dfrac{1}{3^n}=\dfrac{-2}{3.3^n} < 0$ nên $u_n=\dfrac{1}{3^n}$ là dãy số  giảm.\\
		Xét $u_n=\dfrac{1}{2n+1}$. Ta có\\ 
		$u_{n+1}-u_n=\dfrac{1}{2\left(n+1\right)+1}-\dfrac{1}{2n+1}=\dfrac{1}{2n+3}-\dfrac{1}{2n+1}=\dfrac{-2}{\left(2n+3\right)\left(2n+1\right)} < 0$ nên $u_n=\dfrac{1}{2n+1}$ là dãy  số giảm.\\
		Xét $u_n=\dfrac{n+1}{3n+2}$. Ta  có\\
		$u_{n+1}-u_n=\dfrac{n+2}{3n+5}-\dfrac{n+1}{3n+2}=-\dfrac{1}{\left(3n+5\right)\left(3n+2\right)} < 0$  nên $u_n=\dfrac{n+1}{3n+2}$ là dãy số giảm.\\
		Xét $u_n=\dfrac{4n-2}{n+3}$. Ta  có\\
		$u_{n+1}-u_n=\dfrac{4n+2}{n+4}-\dfrac{4n-2}{n+3}=\dfrac{14}{\left(n+4\right)\left(n+3\right)} > 0$ nên $u_n=\dfrac{4n-2}{n+3}$ là dãy  số tăng.
	} 
\end{ex}

% nguồn: Ok_Toan11_TracNghiem_Lan2
\begin{ex}%[1D2H1-4]%[ma:a321b864]
	Trong các dãy số sau, dãy số nào là dãy số giảm 
	 \choice
		{$u_n=\dfrac{n-3}{n+1}$} 
		{$u_n=\dfrac{n}{2}$} 
		{\True $u_n=\dfrac{2}{n^2}$} 
		{$u_n=\dfrac{(-1)^n}{3^n}$} 
	\loigiai{\begin{itemize}
		\item Ta có $u_n = \dfrac{n-3}{n+1}$; $u_{n+1}  = \dfrac{n-2}{n+2}$.\\ Khi đó $u_{n+1}-u_n = \dfrac{n-2}{n+2}-\dfrac{n-3}{n+1} = \dfrac{4}{(n+1) (n+2) } > 0$, $\forall n \in \mathbb{N}$.\\ Vậy($u_n$) là dãy số tăng.
		\item Ta có $u_n = \dfrac{n}{2}$; $u_{n+1}  = \dfrac{n+1}{2}$.\\
		Khi đó $u_{n+1}-u_n =  \dfrac{n+1}{2}-\dfrac{n}{2} = \dfrac{1}{2} > 0$, $\forall n \in \mathbb{N}$.\\
		Vậy($u_n$) là dãy số tăng. 
		\item Ta có $u_n = \dfrac{2}{n^2}$; $u_{n+1} = \dfrac{2}{(n+1)^2}$.\\ 
		Khi đó $\dfrac{u_{n+1}}{u_{n}}  = \dfrac{n^2}{(n+1)^2} < \dfrac{n^2}{n^2} = 1$, $\forall n \in \mathbb{N}^*$.\\
		Vậy($u_n$) là dãy giảm. 
		\item Ta có $u_1 = \dfrac{-1}{3}$; $u_2 = \dfrac{1}{9}$;  $u_3 = \dfrac{-1}{27}$.\\
		Vậy($u_n$) là dãy số không tăng  không giảm.
\end{itemize}}
\end{ex}

% nguồn: Ok_Toan11_TracNghiem_Lan2
\begin{ex}%[1D2H1-4]%[ma:ab9061ac]
	Trong các dãy số dưới đây, dãy số  nào là dãy giảm?
	 \choice
		{$u_n=n^2$} 
		{\True  $u_n=\dfrac{1}{n}-3$}
		 {$u_n=3n$}
		{$u_n=n^3-2$} 
	\loigiai{Xét đáp án $u_n=n^2$, ta có  $u_{n+1}-u_n=(n+1)^2-n^2=2n+1>0$, $\forall n \in \mathbb{N}^*$ nên dãy này là dãy tăng.\\
		Xét đáp án $u_n=\dfrac{1}{n}-3$, ta có $u_{n+1}-u_n=\dfrac{1}{n+1}-\dfrac{1}{n}=\dfrac{-1}{n(n+1) }<0$, $\forall n \in \mathbb{N}^*$  nên dãy này là dãy giảm.\\
		Xét đáp án $u_n=3n$,  ta có $u_{n+1}-u_n=3(n+1) -3n=3>0$, $\forall n \in \mathbb{N}^*$ nên dãy này là dãy tăng.\\
	Xét đáp án $u_n=n^3-2$, ta có $u_{n+1}-u_n=(n+1)^3-n^3>0$,  $\forall n \in \mathbb{N}^*$ nên dãy này là dãy tăng.}
\end{ex}

% nguồn: Ok_Toan11_TracNghiem_Lan2
\begin{ex}%[1D2H1-4]%[ma:4ef53773]
	Trong các  dãy số $(u_n)$ cho bởi số hạng tổng quát $u_n$ sau, dãy số nào là dãy số tăng?
	 \choice
		 {$u_n=\dfrac{1}{2^n}$}
		 {$u_n=\dfrac{1}{n}$}
		 {$u_n=\dfrac{n+5}{3n+1}$}
		 {\True $u_n=\dfrac{2n-1}{n+1}$}
	\loigiai{\begin{itemize}
		\item Vì $2^n$, $n$ là các dãy dương và tăng  nên $\dfrac{1}{2^n}$, $\dfrac{1}{n}$ là các dãy giảm.
		\item Xét dãy số $u_n=\dfrac{n+5}{3n+1}$ $u_n = \dfrac{n+5}{3n+1} \Rightarrow\heva{ &u_1 = \dfrac{3}{2} \\  &u_2 = \dfrac{7}{6}} \Rightarrow u_1 > u_2$. \item Xét dãy số $u_n=\dfrac{2n-1}{n+1}$ $u_n = \dfrac{2n-1}{n+1} = 2 - \dfrac{3}{n+1} \Rightarrow u_{n+1}-u_n = 3\left(\dfrac{1}{n+1} - \dfrac{1}{n+2} \right) > 0.$
\end{itemize}}
\end{ex}

% nguồn: Ok_Toan11_TracNghiem_Lan2
\begin{ex}%[1D2H1-4]%[ma:b09abec2]
	Cho dãy số $\left(u_n \right)$ biết $\heva{&u_1=1 \\& u_{n+1}=au_n+1, \forall n \in \mathbb{N}^*.}$ Tìm tất cả các giá trị của $a$ để $\left(u_n \right)$ tăng. 
	\choice 
		{$a < 0$} 
		{$a  \le 0$}
		{\True $a  > 0$}
		{$a  > 1$}
	\loigiai{Xét hiệu $u_{n+1}=au_n+1 \Rightarrow  u_n=au_{n-1}+1 \Rightarrow u_{n+1}-u_n=a(u_n-u_{n-1})$\\ $\Rightarrow u_3-u_2=a(u_2-u_1) =a^2 \Rightarrow u_4-u_3=a(u_3-u_2) =a^3\cdots\Rightarrow u_{n+1}-u_n=a^n$.\\
	Để dãy số $\left(u_n \right)$ tăng thì $a^n.0$ suy ra $a > 0$.} 
\end{ex}

% nguồn: Ok_Toan11_TracNghiem_Lan2
\begin{ex}%[1D2H1-4]%[ma:6f15ce0d]
	Trong các dãy số $(u_n)$ cho bởi số hạng tổng quát  $u_n$ sau, dãy số nào tăng?
	\choice 
		{$u_n=\dfrac{\sin  n}{n}$}
		{$u_n=\dfrac{\sqrt{n^2+1}}{2n+1}$} 
		{$u_n=\dfrac{3^n}{n^2}$} 
		 {\True $u_n=4n^3-3n^2+1$}
	\loigiai{Với $n \in(k2\pi;\pi+k2\pi)$, $k \in  \mathbb{N} \Rightarrow \sin n > 0 \Rightarrow \dfrac{\sin n}{n} > 0$.\\
		 Với $n \in(\pi+k2\pi;2\pi+k2\pi)$, $k \in \mathbb{N} \Rightarrow \sin n < 0 \Rightarrow \dfrac{\sin n}{n} < 0$.\\
		Ta có $u_n = \dfrac{\sqrt{n^2+1}}{2n+1} =  \sqrt{\dfrac{n^2+1}{(2n+1)^2}}$. Xét dãy $(v_n)$ với $v_n = \dfrac{n^2+1}{(2n+1)^2}$.\\ $v_{n+1}-v_n = \dfrac{n^2+2n+2}{4n^2+12n+9} - \dfrac{n^2+1}{4n^2+4n+1} = \dfrac{4n^2-2n-7}{(2n+3)^2(2n+1)^2}$.\\
		Do $v_{n+1}-v_n$ vừa nhận giá trị âm lẫn dương nên  dãy số $(v_n)$ không tăng, không giảm.\\
		$u_{n+1}-u_n  = \dfrac{3\cdot3^n}{(n+1)^2} - \dfrac{3^n}{n^2} = \dfrac{3^n(3n^2-2n-1) }{(n+1)^2n^2}$.\\
	Do  $u_{n+1}-u_n$ nhận giá trị âm lẫn dương nên dãy đã cho không tăng, không giảm.}
\end{ex}

% nguồn: Ok_Toan11_TracNghiem_Lan2
\begin{ex}%[1D2H1-4]%[ma:5ec6444d]
	Trong  các dãy số sau, dãy số nào là dãy số giảm?
	\choice 
		{$u_n=n^2$} 
		 {$u_n=2n$}
		 {$u_n=n^3-1$}
		 {\True $u_n=-2n+1$}
	\loigiai{ 
	Xét dãy $u_n=n^2$.  Ta có $u_{n+1}-u_n=2n +1> 0, \forall n\in \mathbb{N}^{*}$, suy ra Dãy số tăng.\\
	Xét  dãy $u_n=2n$. Ta có $u_{n+1}-u_n=2> 0, \forall n\in \mathbb{N}^{*}$, suy ra Dãy số tăng.\\
	Xét dãy $u_n=n^3-1$. Ta có $u_{n+1}-u_n=3n^2+3n+1> 0, \forall n\in  \mathbb{N}^{*}$, suy ra Dãy số tăng.\\
	Ta có $u_{n+1}=-2(n+1)+1=-2n-1$, $u_{n+1}-u_n=-2<  0$.\\
	 Vậy $u_n=-2n+1$ là hàm số giảm.
	} 
\end{ex}

% nguồn: Ok_Toan11_TracNghiem_Lan2
\begin{ex}%[1D2H1-4]%[ma:e59100ec]
	Cho dãy số $(u_n)$ biết $\heva{  &u_1=1 \\ &u_n=\dfrac{1}{3}u_{n-1}+\dfrac{5}{3}}$. Mệnh đề nào sau đây đúng?
	 \choice
		{\True Dãy số  tăng}
		 {Dãy số giảm}
		{Dãy số không tăng, không  giảm}
		{Dãy số vừa  tăng vừa giảm}
	 \loigiai{
		 Ta có $u_1 < u_2 < u_3$. Dự đoán dãy số đã cho tăng, ta chứng minh bằng quy nạp.\\
		Từ giả thiết thì $u_n > 0$, $\forall  n \in \mathbb{N}^*$.\\
		Giả sử $u_k > u_{k-1}$, $k \ge 2$. Ta chứng minh $u_{k+1}  > u_k$.\\
		Thật vậy, $u_{k+1}-u_k = \dfrac{1}{3}\left(u_{k-1}-u_k\right)  > 0 \Leftrightarrow u_{k+1} > u_k$.\\
	Vậy dãy đã  cho là dãy tăng.}
\end{ex}

% nguồn: Ok_Toan11_TracNghiem_Lan2
\begin{ex}%[1D2H1-4]%[ma:9a03ed6b]
	Dãy  số nào sau đây là dãy số giảm?
	 \choice
		{\True  $u_n=\dfrac{5-3n}{2n+3}$, $(n \in \mathbb{N}^*)$}
		{$u_n=\dfrac{n-5}{4n+1}$, $(n \in  \mathbb{N}^*)$}
		{$u_n=2n^3+3$, $(n \in \mathbb{N}^*)$} 
		{$u_n=\cos(2n+1)$, $(n \in  \mathbb{N}^*)$}
	\loigiai{Xét $u_n = \dfrac{5-3n}{2n+3}$, $(n \in \mathbb{N}^*)$, ta có \begin{eqnarray*}
		u_{n+1}-u_n &=& \dfrac{5-3(n+1)  }{2(n+1) +3} - \dfrac{5-3n}{2n+3}\\
		&=& \dfrac{2-3n}{2n+5} - \dfrac{5-3n}{2n+3}\\ 
		&=& \dfrac{(2-3n) (2n+3) -(2n+5) (5-3n) }{(2n+5)  (2n+3) }\\
		&=&  \dfrac{4n-6n^2+6-9n-10n+6n^2-25+15n}{(2n+5) (2n+3) }\\
		&=& \dfrac{-19}{(2n+5) (2n+3) }  < 0, \forall n \in \mathbb{N}^*.
	\end{eqnarray*}
Vậy $u_n = \dfrac{5-3n}{2n+3}$, $(n \in \mathbb{N}^*)$ là  dãy giảm.}
\end{ex}

% nguồn: Ok_Toan11_TracNghiem_Lan2
\begin{ex}%[1D2H1-4]%[ma:a7f9a7ef]
	 Dãy số nào sau đây là dãy số giảm?
	 \choice
		{\True $u_n=\dfrac{5-3n}{2n+3}$,  $n \in \mathbb{N}^*$}
		{$u_n=\dfrac{n-5}{4n+1}$, $n \in \mathbb{N}^*$} 
		{$u_n=2n^2+3$, $n \in \mathbb{N}^*$} 
		 {$u_n=\cos(2n+1)$, $n \in \mathbb{N}^*$}
	\loigiai{Với dãy $u_n = \dfrac{5-3n}{2n+3}$.\\ Ta có \begin{eqnarray*}
		u_{n+1} - u_n  &=& \dfrac{5-3(n+1) }{2(n+1) +3} - \dfrac{5-3n}{2n+3} \\
		&=& \dfrac{2-3n}{2n+5} - \dfrac{5-3n}{2n+3}\\ 
		&=& \dfrac{(2-3n)  (2n+3) -(5-3n) (2n+5) }{(2n+5) (2n+3) }\\
		&=& \dfrac{-19}{(2n+3) (2n+5) } < 0, \forall  n \in \mathbb{N}^*.
	\end{eqnarray*}
Vậy $(u_n)$ là  dãy giảm.}
\end{ex}

% nguồn: Ok_Toan11_TracNghiem_Lan2
\begin{ex}%[1D2H1-4]%[ma:7d961bd5]
	Cho dãy số($u_n$) biết $u_n=\sqrt{3n+2}-\sqrt{3n+1}$. Mệnh  đề nào sau đây đúng?
	\choice 
		{Dãy số tăng} 
		{\True Dãy số giảm} 
		{Dãy số không  tăng, không giảm}
		{Dãy số vừa tăng vừa giảm} 
	\loigiai{Ta có $u_n = \sqrt{3n+2}  - \sqrt{3n+1} = \dfrac{1}{\sqrt{3n+2} + \sqrt{3n+1}}$.\\
		Khi đó \begin{eqnarray*}
		u_{n+1}-u_n &=& \dfrac{1}{\sqrt{3n+5} + \sqrt{3n+4}}  - \dfrac{1}{\sqrt{3n+2} + \sqrt{3n+1}}\\
		&=& \dfrac{(\sqrt{3n+2}-\sqrt{3n+5}) +(\sqrt{3n+1}-\sqrt{3n+4})  }{(\sqrt{3n+5}+\sqrt{3n+4}) (\sqrt{3n+2}+\sqrt{3n+1}) } < 0, \forall n \in \mathbb{N}^*.
\end{eqnarray*}}
\end{ex}

% nguồn: Ok_Toan11_TracNghiem_Lan2
\begin{ex}%[1D2H1-4]%[ma:fb0a4ab7]
	Cho dãy số $(u_n)$ biết $u_n=\dfrac{1}{3n+2}$. Mệnh đề nào sau đây là  đúng?
	 \choice
		 {Dãy số tăng}
		{\True  Dãy số giảm}
		{Dãy số không tăng, không giảm} 
		{Dãy số vừa tăng vừa giảm} 
	 \loigiai{
		Ta có $u_{n+1}-u_n=\dfrac{1}{3(n+1)+2}-\dfrac{1}{3n+2}=\dfrac{1}{3n+5}-\dfrac{1}{3n+2}=-\dfrac{3}{(3n+5)(3n+2)}  < 0$, $\forall n\in\mathbb{N}^*$.\\
		Vậy $u_{n+1}-u_n < 0\Leftrightarrow u_{n+1} < u_n$, $\forall n  \in \mathbb{N}^*$.
	} 
\end{ex}

% nguồn: Ok_Toan11_TracNghiem_Lan2
\begin{ex}%[1D2H1-4]%[ma:407a4e1e]
	Trong các dãy  số $(u_n)$ cho bởi số hạng tổng quát $u_n$ sau, dãy số nào không tăng, không giảm?
	 \choice
		 {$u_n=n+\dfrac{1}{n}$}
		 {$u_n=5^n+3n$}
		{$u_n=-3^n$} 
		{\True $u_n=(-3)^n  \cdot \sqrt{n^2+1}$}
	\loigiai{ 
		Xét  dãy số có số hạng tổng quát $u_n=(-3)^n \cdot \sqrt{n^2+1}$.\\
		Dãy số này không tăng, không giảm vì các số  hạng đan dấu.
	 }
\end{ex}

% nguồn: Ok_Toan11_TracNghiem_Lan2
\begin{ex}%[1D2H1-4]%[ma:1b739c4c]
	Trong các dãy số $(u_n)$ cho  bởi số hạng tổng quát $u_n$ sau, dãy số nào là dãy số tăng?
	 \choice
		 {$u_n=\dfrac{2}{3^n}$}
		 {$u_n=\dfrac{3}{n}$}
		 {\True $u_n=2^n$}
		{$u_n=(-2)^n$} 
	\loigiai{Xét dãy số $u_n=2^n$ $u_n =  2^n \Rightarrow u_{n+1} - u_n = 2^{n+1} - 2^n = 2^n > 0$.\\
	Vậy dãy số $(u_n)$  với $u_n=2^n$ là dãy số tăng.}
\end{ex}

% nguồn: Ok_Toan11_TracNghiem_Lan2
\begin{ex}%[1D2H1-4]%[ma:0a8081cf]
	Trong các dãy số $(u_n)$ cho bởi số hạng tổng quát $u_n$  sau, dãy số nào tăng?
	 \choice
		 {$u_n=\dfrac{1}{3^n}$}
		 {$u_n=\dfrac{1}{2n+1}$}
		{$u_n=\dfrac{n+1}{3n+2}$} 
		 {\True $u_n=\dfrac{4n-2}{n+3}$}
	 \loigiai{
		Ta  có
		\begin{itemize}
			\item $u_{n+1}-u_n=\dfrac{1}{3^{n+1}}-\dfrac{1}{3^n}=\dfrac{1}{3\cdot 3^n}-\dfrac{1}{3^n}=\dfrac{-2}{3\cdot 3^n} < 0$, $\forall  n\in\mathbb{N}^*$.
			\item$u_{n+1}-u_n=\dfrac{1}{2(n+1)+1}-\dfrac{1}{2n+1}=\dfrac{1}{2n+3}-\dfrac{1}{2n+1}=\dfrac{-2}{(2n+3)(2n+1)} < 0$, $\forall n\in\mathbb{N}^*$. 
			\item $u_{n+1}-u_n=\dfrac{n+2}{3n+5}-\dfrac{n+1}{3n+2}=-\dfrac{1}{(3n+5)(3n+2)} < 0$, $\forall n\in\mathbb{N}^*$. 
			\item $u_{n+1}-u_n=\dfrac{4n+2}{n+4}-\dfrac{4n-2}{n+3}=\dfrac{14}{(n+4)(n+3)}  > 0$, $\forall n\in\mathbb{N}^*$.
		\end{itemize}
		Vậy dãy số $(u_n)$ cho bởi số hạng tổng  quát $u_n=\dfrac{4n-2}{n+3}$ là dãy số tăng.
	} 
\end{ex}

% nguồn: Ok_Toan11_TracNghiem_Lan2
\begin{ex}%[1D2H1-4]%[ma:5aade844]
	Cho dãy số $(u_n)$ biết $u_n=\dfrac{2n^2-n-1}{n+2}$. Mệnh  đề nào sau đây đúng?
	 \choice
		 {\True Dãy số tăng}
		 {Dãy số giảm}
		{Dãy số không  tăng, không giảm}
		 {Có số hạng âm}
	\loigiai{ 
		Ta có $u_{n+1}-u_n=\dfrac{2n^2+3n}{n+3}-\dfrac{2n^2-n-1}{n+2}=\dfrac{2n^2+10n+3}{\left(n+3\right)\left(n+2\right)}  > 0$, $ \forall n\in \mathbb{N}^{*}$.\\
		Vậy dãy số đã cho là dãy tăng. 
	 }
\end{ex}

% nguồn: Ok_Toan11_TracNghiem_Lan2
\begin{ex}%[1D2H1-4]%[ma:fddd006b]
	 Cho dãy số $(u_n)$ biết $u_n=2n^2+3n+1$. Mệnh đề nào sau đây đúng?
	 \choice
		{Dãy số vừa tăng vừa  giảm}
		{Dãy số  giảm}
		{Dãy số không tăng, không giảm} 
		{\True Dãy số  tăng}
	\loigiai{ 
		Ta có $u_{n+1}-u_n=2\left(n+1\right)^2+3\left(n+1\right)+1-2n^2-3n-1=4n+5>  0$, $ \forall n\in \mathbb{N}^{*}$.\\
		Vậy $u_{n+1}-u_n > 0\Leftrightarrow u_{n+1} > u_n$, $  \forall n\in \mathbb{N}^{*}$.
	} 
\end{ex}

% nguồn: Ok_Toan11_TracNghiem_Lan2
\begin{ex}%[1D2H1-4]%[ma:908fe078]
	Trong  các dãy số $(u_n)$ cho bởi số hạng tổng quát $u_n$ sau, dãy số nào giảm?
	 \choice
		 {$u_n=\left(\dfrac{4}{3}\right)^n$}
		 {$u_n=(-1)^n\left(5^n-1\right)$}
		{\True  $u_n=-3^n$}
		 {$u_n=\sqrt{n+4}$}
	\loigiai{ 
		Ta có $u_{n+1}-u_n=\left(\dfrac{4}{3}\right)^{n+1}-\left(\dfrac{4}{3}\right)^n=\dfrac{4}{3} \cdot\left(\dfrac{4}{3}\right)^n-\left(\dfrac{4}{3}\right)^n=\dfrac{1}{3}  \cdot\left(\dfrac{4}{3}\right)^n > 0$, $\forall n\in\mathbb{N}^*$.\\
		Dãy $(u_n)$ với $u_n=(-1)^n\left(5^n-1\right)$ có các số hạng đan dấu nên  dãy không tăng, không giảm.\\
		$ u_{n+1}-u_n=-3^{n+1}+3^n=-3\cdot 3^n+3^n=-2\cdot 3^n < 0  $, $\forall n\in\mathbb{N}^*$.\\
		$u_{n+1}-u_n=\sqrt{n+5}-\sqrt{n+4}=\dfrac{1}{\sqrt{n+5}+\sqrt{n+4}} > 0$,  $\forall n\in\mathbb{N}^*$.
	 }
\end{ex}

% nguồn: Ok_Toan11_TracNghiem_Lan2
\begin{ex}%[1D2H1-4]%[ma:750a0f38]
	Trong các dãy số $(u_n)$ cho  bởi số hạng tổng quát sau, dãy số nào là dãy số tăng?
	 \choice
		 {$u_n = 1 - 2n$}
		{$u_n  = 3$}
		{\True  $u_n = 2n$}
		{$u_n = (-1)^n$} 
	\loigiai{ 
		Xét dãy số $(u_n) \colon u_n =  2n$. Ta có
		\[ 
		u_{n+1} - u_n = 2(n+1) - 2n = 2 > 0, \  \forall n \in \mathbb{N}^*.
		\] 
		Do đó dãy $u_n = 2n$ là  dãy số tăng.
	 }
\end{ex}

% nguồn: Ok_Toan11_TracNghiem_Lan2
\begin{ex}%[1D2H1-4]%[ma:0e7c602c]
	Cho dãy số $(u_n)$ biết $u_n=n^2-400n$. Mệnh  đề nào sau đây đúng?
	\choice 
		{Dãy số tăng} 
		{Dãy  số giảm}
		{\True  Dãy số không tăng, không giảm}
		{Mọi số  hạng đều âm}
	 \loigiai{
		 Ta có $u_{n+1}-u_n=\left(n+1\right)^2-400\left(n+1\right)-n^2+400n=2n-399$.\\
		Do $2n-399> 0$ khi $n > \dfrac{399}{2}$ và $2n-399< 0$  khi $n < \dfrac{399}{2}$.\\
		 Vậy dãy số đã cho không tăng, không giảm.
	 }
\end{ex}

% nguồn: Ok_Toan11_TracNghiem_Lan2
\begin{ex}%[1D2H1-4]%[ma:2aeae3c1]
	Cho dãy số($u_n$) biết $u_n=n-\sqrt{n^2+1}$. Mệnh đề nào sau đây  đúng?
	\choice 
		{\True Dãy số tăng} 
		{Dãy số giảm} 
		 {Dãy số không tăng, không giảm}
		{Các số hạng  đều dương}
	\loigiai{Ta có $u_n = n - \sqrt{n^2+1} =  \dfrac{-1}{n+\sqrt{n^2+1}}$.\\
		Khi đó\begin{eqnarray*}
		u_{n+1}-u_n &=&  \dfrac{-1}{n+1+\sqrt{(n+1)^2+1}} - \dfrac{-1}{n+\sqrt{n^2+1}}\\
		 &=& \dfrac{1}{n+\sqrt{n^2+1}} - \dfrac{1}{n+1+\sqrt{(n+1)^2+1}}\\
		&=& \dfrac{1+(\sqrt{(n+1)^2+1}-\sqrt{n^2+1}) }{(n+1+\sqrt{(n+1)^2+1})  (n+\sqrt{n^2+1}) } > 0, \forall n \in \mathbb{N}^*.
\end{eqnarray*} Vậy dãy số đã cho là dãy tăng.}
\end{ex}

% nguồn: Ok_Toan11_TracNghiem_Lan2
\begin{ex}%[1D2H1-4]%[ma:97bdda92]
	Trong các dãy số $\left(u_n \right)$ cho bởi số  hạng tổng quát $u_n$ sau, dãy số nào không tăng, không giảm?
	 \choice
		{$u_n=n+\dfrac{1}{n}$} 
		 {$u_n=5^n+3n$}
		{$u_n=-3^n$} 
		{\True $u_n=\left(-3\right)^n.\sqrt{n^2+1}$} 
	\loigiai{ 
		Xét $u_n=n+\dfrac{1}{n}$ có $  u_{n+1}-u_n= \dfrac{n^2+n-1}{n^2+n}>0 $ $ \forall n\in \mathbb{N}^{*}$. Do đó dãy số tăng.\\
		Xét $u_n=5^n+3n$ có $  u_{n+1}-u_n= 4\cdot 5^{n}+3>0 $ $ \forall n\in \mathbb{N}^{*}$. Do đó dãy số tăng.\\
		 Xét $u_n=-3^n$ có $ u_{n+1}-u_n= -2\cdot 3^n <0 $ $ \forall n\in \mathbb{N}^{*}$. Do đó dãy số giảm.\\
		Dãy $u_n=\left(-3\right)^n\cdot\sqrt{n^2+1}$ có $ u_1 = -3\sqrt{2}, u_2 = 9\sqrt{5}, u_3 = -27\sqrt{10} $ nên dãy không tăng, không giảm. 
	} 
\end{ex}

% nguồn: Ok_Toan11_TracNghiem_Lan2
\begin{ex}%[1D2H1-4]%[ma:0e1133df]
	Trong các dãy số sau, dãy số nào  là dãy số giảm?
	\choice 
		{\True $u_n=\dfrac{2n+1}{n-1}$} 
		{$u_n=n^3-1$} 
		{$u_n=n^2$} 
		{$u_n=2n$} 
	\loigiai{Xét dãy số $u_n=\dfrac{2n+1}{n-1}$. Với mọi  $n \in \mathbb{N}$, $n > 1$.\\
		Ta có \begin{eqnarray*}
		u_{n+1}-u_n &=& \dfrac{2(n+1) +1}{(n+1) -1}  - \dfrac{2n+1}{n-1}\\
		 &=& \dfrac{2n+3}{n} - \dfrac{2n+1}{n-1}\\
		&=& \dfrac{(2n+3) (n-1) -n(2n+1) }{n(n-1) }  \\
		&=&  \dfrac{(2n+3) (n-1) -n(2n+1) }{n(n-1) }\\
		&=& \dfrac{-3}{n(n-1) } < 0, \forall n \in \mathbb{N}, n  > 1.
\end{eqnarray*} Suy ra dãy số giảm.}
\end{ex}

% nguồn: Ok_Toan11_TracNghiem_Lan2
\begin{ex}%[1D2H1-4]%[ma:22dcf9d3]
	Cho dãy số $(u_n)$ biết $u_n=2n^2+3n+1$.  Mệnh đề nào sau đây là đúng?
	\choice 
		{\True Dãy số tăng} 
		 {Dãy số giảm}
		{Dãy số  không tăng, không giảm}
		{Dãy  số vừa tăng vừa giảm}
	 \loigiai{
		Ta có $u_{n+1}-u_n=2(n+1)^2+3(n+1)+1-2n^2-3n-1=4n+5> 0, \forall n  \in \mathbb{N}^*$.\\
		Vậy $u_{n+1}-u_n  < 0\Leftrightarrow u_{n+1} < u_n, \forall n \in \mathbb{N}^*$
	} 
\end{ex}

% nguồn: Ok_Toan11_TracNghiem_Lan2
\begin{ex}%[1D2H1-4]%[ma:c839554f]
	Trong các dãy số $\left(u_n \right)$ cho bởi số hạng  tổng quát $u_n$ sau, dãy số nào giảm?
	 \choice
		{$u_n=\left(\dfrac{4}{3}  \right)^n$}
		{$u_n=\left(-1\right)^n  \left(5^n-1\right)$}
		{\True  $u_n=-3^n$}
		{$u_n=\sqrt{n+4}$} 
	 \loigiai{
		Dãy số $u_n=\left(\dfrac{4}{3} \right)^n$ có $u_{n+1}-u_n=\left(\dfrac{4}{3} \right)^{n+1}-\left(\dfrac{4}{3} \right)^n=\dfrac{4}{3}. \left(\dfrac{4}{3} \right)^n-\left(\dfrac{4}{3} \right)^n=\dfrac{1}{3}. \left(\dfrac{4}{3} \right)^n > 0$. Do  đó $u_n=\left(\dfrac{4}{3} \right)^n$ là dãy số tăng.\\
		Dãy số $\left(u_n \right)$ với $u_n=\left(-1\right)^n \left(5^n-1\right)$ có $ u_1 =-4, u_2 = 24, u_ 3 = -124 $  nên $u_n=\left(-1\right)^n \left(5^n-1\right)$ không giảm.\\
		Dãy số $u_n=\sqrt{n+4}$ có  $u_{n+1}-u_n=\sqrt{n+5}-\sqrt{n+4}=\dfrac{1}{\sqrt{n+5}+\sqrt{n+4}} > 0$. Do đó $u_n=\sqrt{n+4}$ là dãy số tăng.\\
		Dãy số $u_n=-3^n$ có $u_{n+1}-u_n=-3^{n+1}+3^n=-3.3^n+3^n=-2.3^n < 0$ nên là  dãy số giảm.
	} 
\end{ex}

% nguồn: Ok_Toan11_TracNghiem_Lan2
\begin{ex}%[1D2H1-4]%[ma:78714c55]
	Trong các dãy số  ($u_n$) cho bởi số hạng tổng quát $u_n$ sau, dãy số nào là dãy số giảm?
	\choice 
		 {\True $u_n=\dfrac{1}{2^n}$}
		{$u_n=\dfrac{3n-1}{n+1}$} 
		{$u_n=n^2$} 
		{$u_n=\sqrt{n+2}$} 
	\loigiai{Xét  dãy số $u_n=\dfrac{1}{2^n}$. Ta có $u_n = \dfrac{1}{2^n} < \dfrac{1}{2^{n+1}} = u_{n+1}$, $\forall n \in \mathbb{N}^*$.}
\end{ex}

% nguồn: Ok_Toan11_TracNghiem_Lan2
\begin{ex}%[1D2H1-4]%[ma:b62e45f0]
	 Cho dãy số $(u_n)$ biết $\heva{ &u_1=1 \\& u_{n+1}=\sqrt{u_n^2+3}, \forall n \in \mathbb{N}^*.}$ Mệnh đề nào sau đây đúng?
	 \choice
		{\True Dãy số tăng} 
		{Dãy số giảm} 
		{Dãy số không tăng,  không giảm}
		{Dãy số vừa tăng vừa  giảm}
	\loigiai{Ta có $0 <  u_1 < u_2 < u_3$. Dự đoán dãy số đã cho tăng, ta chứng minh bằng quy nạp.\\
		Từ giả thiết thì $u_n>0$, $\forall n\in\mathbb{N}^*$. Giả  sử $u_k>u_{k-1}$, $k\geq 2$. Ta chứng minh $u_{k+1}>u_k$.\\
		Thật vậy \begin{eqnarray*}
		u_{k+1}-u_k&=&\sqrt{u_k^2+3}-\sqrt{u_{k-1}^2+3}\\ 
		 &=&\dfrac{\left(u_k-u_{k-1}\right)\left(u_k+u_{k-1}\right)}{\sqrt{u_k^2+3}+\sqrt{u_{k-1}^2+3}}>0\\
		&\Rightarrow&  u_{k+1}>u_k.
	\end{eqnarray*}
	Vậy dãy số đã cho là dãy  tăng.
 }
\end{ex}

% nguồn: Ok_Toan11_TracNghiem_Lan2
\begin{ex}%[1D2H1-4]%[ma:58a4eaf7]
	Cho dãy số $(u_n)$ biết $u_n=\sqrt{5n+2}$. Mệnh đề nào sau đây  là đúng?
	\choice 
		{\True  Dãy số tăng}
		{Dãy  số giảm}
		 {Dãy số không tăng, không giảm}
		{Dãy số vừa  tăng vừa giảm}
	 \loigiai{
		Ta có $u_{n+1}-u_n=\sqrt{5(n+1)+2}-\sqrt{5n+2}=\sqrt{5n+7}-\sqrt{5n+2} > 0\Leftrightarrow u_{n+1} > u_n$, $\forall  n\in\mathbb{N}^*$.
	} 
\end{ex}

% nguồn: Ok_Toan11_TracNghiem_Lan2
\begin{ex}%[1D2H1-4]%[ma:42435af4]
	Cho dãy số $(u_n)$ biết $u_n=\dfrac{10}{3^n}$. Mệnh  đề nào sau đây là đúng?
	\choice 
		{Dãy số tăng} 
		{\True Dãy số  giảm}
		{Dãy số không tăng, không  giảm}
		{Dãy số vừa tăng vừa giảm} 
	 \loigiai{
		Ta có $u_{n+1}-u_n=\dfrac{10}{3^{n+1}}-\dfrac{10}{3^n}=\dfrac{10}{3\cdot3^n}-\dfrac{10}{3^n}=\dfrac{-20}{3\cdot3^n} < 0$, $\forall n \in  \mathbb{N}^*$.\\
		Vậy $u_{n+1}-u_n < 0\Leftrightarrow u_{n+1} < u_n, \forall n  \in \mathbb{N}^*$.
	 }
\end{ex}

% nguồn: Ok_Toan11_TracNghiem_Lan2
\begin{ex}%[1D2H1-4]%[ma:9827367c]
	Cho dãy số $(u_n)$ biết $u_n=(-1)^n\left(n^2+1\right)$. Mệnh đề nào sau đây là  đúng?
	 \choice
		 {Dãy số tăng}
		{Dãy số giảm} 
		{\True Dãy số không  tăng, không giảm}
		{Dãy số là  dãy hữu hạn}
	 \loigiai{
		Ba số hạng đầu  tiên của dãy số lần lượt là $-2$; $5$; $-10$. Từ đó ta có thể kết luận được dãy số đã cho là dãy không tăng, không giảm.
	} 
\end{ex}

% nguồn: Ok_Toan11_TracNghiem_Lan2
\begin{ex}%[1D2H1-4]%[ma:556bf1c4]
	 Cho dãy số $(u_n)$ biết $u_n=5^n-4^n$. Mệnh đề nào sau đây là đúng?
	 \choice
		{\True $(u_n)$ là  dãy số tăng}
		 {$(u_n)$ là dãy số giảm}
		{$(u_n)$ là dãy số không tăng, không giảm} 
		{$(u_n)$ là dãy số có  số hạng thứ $100$ bé hơn $1$}
	\loigiai{ 
		 Ta có $u_{n+1}-u_n=5^{n+1}-4^{n+1}-5^n+4^n=4\left(5^n-4^n\right) > 0, \forall n \in \mathbb{N}^*$.\\
		Vậy $u_{n+1}-u_n > 0\Leftrightarrow  u_{n+1} > u_n, \forall n \in \mathbb{N}^*$.
	 }
\end{ex}

% nguồn: Ok_Toan11_TracNghiem_Lan2
\begin{ex}%[1D2H1-4]%[ma:7781f786]
	Trong các dãy số ($u_n$) cho bởi số  hạng tổng quát $u_n$ sau, dãy số nào là dãy số tăng?
	\choice 
		{$u_n=\dfrac{1}{2^n}$} 
		{\True $u_n=\dfrac{3n-1}{n+1}$} 
		 {$u_n=1-n^2$}
		 {$u_n=\dfrac{1}{\sqrt{n+2}}$}
	\loigiai{Xét dãy số $u_n=\dfrac{3n-1}{n+1}$. Ta  có $u_{n+1}-u_n=\dfrac{3n+2}{n+2}-\dfrac{3n-1}{n+1}=\dfrac{4}{(n+2)(n+1)}>0$, $\forall n\in\mathbb{N}^*$.\\
	Vậy dãy số $u_n=\dfrac{3n-1}{n+1}$ là dãy số tăng.} 
\end{ex}

% nguồn: Ok_Toan11_TracNghiem_Lan2
\begin{ex}%[1D2H1-4]%[ma:38537ad0]
	Cho  dãy số $(u_n)$ biết $u_n=\sqrt{5n+2}$. Mệnh đề nào sau đây đúng?
	\choice 
		{\True Dãy  số tăng}
		 {Dãy số giảm}
		{Dãy số không tăng, không giảm} 
		 {Dãy số vừa tăng vừa giảm}
	 \loigiai{
		Ta có $u_{n+1}-u_n=\sqrt{5\left(n+1\right)+2}-\sqrt{5n+2}=\sqrt{5n+7}-\sqrt{5n+2}  = \dfrac{5}{\sqrt{5n+7}+\sqrt{5n+2}} > 0$.\\
		Do đó $ u_{n+1} > u_n$.\\ 
		Vậy $(u_n)$ là  dãy số tăng.
	 }
\end{ex}

% nguồn: Ok_Toan11_TracNghiem_Lan2
\begin{ex}%[1D2H1-4]%[ma:227f17cd]
	Cho dãy số($u_n$) biết $u_n=\dfrac{2n^2-n-1}{n+2}$. Mệnh đề  nào sau đây đúng?
	 \choice
		 {\True Dãy số tăng}
		{Dãy  số giảm}
		{Dãy số không tăng, không giảm} 
		{Có số hạng  âm}
	\loigiai{Ta có \begin{eqnarray*}
		u_{n+1}-u_n &=& \dfrac{2(n+1)^2-(n+1) -1}{(n+1) +2} -  \dfrac{2n^2-n-1}{n+2}\\
		&=& \dfrac{2n^2+3n}{n+3} - \dfrac{2n^2-n-1}{n+2} 
		\\&=& \dfrac{2n^2+10n+3}{(n+3) (n+2)  } > 0, \forall n \in \mathbb{N}^*.
	\end{eqnarray*}
 Vậy dãy số đã cho là dãy tăng.}
\end{ex}

% nguồn: Ok_Toan11_TracNghiem_Lan2
\begin{ex}%[1D2H1-4]%[ma:85a42fc2]
	Cho dãy số $(u_n)$ biết $u_n=n^2-400n$. Mệnh đề nào sau đây là  đúng?
	 \choice
		{Dãy  số tăng}
		{Dãy số giảm} 
		{\True Dãy số không  tăng, không giảm}
		{Mọi số hạng  đều âm}
	 \loigiai{
		Ta có $u_{n+1}-u_n=(n+1)^2-400(n+1)-n^2+400n=2n-399$.\\ 
		Do $2n-399> 0$ khi $n > \dfrac{399}{2}$ và $2n-399< 0$  khi $n < \dfrac{399}{2}$.\\
		 Khi đó ta có $u_{199}=-39999$, $u_{200}=-40000$, $u_{201}=-39999$.\\
		Vậy dãy số đã cho không tăng,  không giảm.
	 }
\end{ex}

% nguồn: Ok_Toan11_TracNghiem_Lan2
\begin{ex}%[1D2H1-4]%[ma:526aeac2]
	Trong các dãy số sau, dãy  số nào là dãy số giảm?
	\choice 
		{\True $u_n=\dfrac{3}{n^2}$} 
		{$u_n=\dfrac{n-3}{n+1}$} 
		{$u_n=\dfrac{n}{2}$} 
		 {$u_n=\dfrac{(-1)^n}{3^n}$}
	\loigiai{\begin{itemize}
		\item Ta có $u_n = \dfrac{3}{n^2}$, $u_{n+1} = \dfrac{3}{(n+1)^2}$.\\ $\dfrac{u_{n+1}}{u_n} = \dfrac{n^2}{(n+1)^2} <\dfrac{n^2}{n^2}=  1$, $\forall n \in \mathbb{N}^*$.\\ Vậy $(u_n)$ là dãy giảm.
		\item Ta có  $u_n = \dfrac{n-3}{n+1}$; $u_{n+1} = \dfrac{n-2}{n+2}$.\\
		Khi đó  $u_{n+1} - u_n = \dfrac{n-2}{n+2} - \dfrac{n-3}{n+1} = \dfrac{4}{(n+1) (n+2) } > 0$, $\forall n \in \mathbb{N^*}$.\\
		Vậy $(u_n)$ là  dãy số tăng.
		\item Ta có $u_n = \dfrac{n}{2}$; $u_{n+1}  = \dfrac{n+1}{2}$.\\
		Khi đó $u_{n+1} - u_n = \dfrac{n+1}{2} - \dfrac{n}{2}  = \dfrac{1}{2} > 0$, $\forall n \in \mathbb{N^*}$.\\
		 Vậy $(u_n)$ là dãy số tăng.
		\item Ta có $u_1 = \dfrac{-1}{3}$; $u_2 = \dfrac{1}{9}$; $u_3 =  \dfrac{-1}{27}$.\\
		Vậy $(u_n)$ là dãy số không tăng không  giảm.
\end{itemize}}
\end{ex}

% nguồn: Ok_TracNghiemLop11_Lan1
\begin{ex}%[1D2H1-4]%[ma:95449f9f]
	Cho dãy số $\left(u_n\right)$ với $u_n=-7n+4$. Khẳng định nào  dưới đây \textbf{sai}?
	\choice 
		{Dãy số  bị chặn trên}
		{Dãy  số không bị chặn}
		{Dãy số giảm} 
		{\True Dãy số tăng} 
	\loigiai{ 
		\begin{itemize}
			\item Ta  có $u_n=-7n+4\Rightarrow{u_{n+1}}-u_n=\left[-7(n+1)+4\right]-\left[-7n+4\right]=-7 < 0,\forall n.$ \\
			Suy ra, $\left(u_n\right)\,$ là  dãy số giảm.
			\item Mà $u_n=-7n+4\le -7\cdot1+4=-3, \forall  n\in\mathbb{N}^*$ nên dãy số bị chặn trên.
			\item Mặt khác, khi $n$ càng lớn thì $u_n$ càng nhỏ, nên dãy số không bị chặn dưới, do đó dãy số không  bị chặn.
		\end{itemize}
	} 
\end{ex}

% nguồn: Ok_TracNghiemLop11_Lan1
\begin{ex}%[1D2H1-4]%[ma:b8d8796d]
	Trong các dãy số  $\left(u_{n} \right)$ được cho bởi số hạng tổng quát sau đây, dãy số nào là dãy số giảm?
	\choice 
		{$u_{n}  =n^{2} ,\, \forall n\in \mathbb{N}^{*} $}
		{$u_{n} =\sqrt{n+1} ,\, \forall n\in  \mathbb{N}^{*} $}
		{$u_{n} =\frac{n^{2}  +1}{n} ,\, \forall n\in \mathbb{N}^{*}$}
		{\True $u_{n} =\frac{1}{2^{n} } ,\, \forall  n\in \mathbb{N}^{*}$}
	\loigiai{ 
		Xét dãy số $u_{n} =\dfrac{1}{2^{n}  } ,\, \forall n\in \mathbb{N}^{*}$.\\
		Có $u_{n} >u_{n+1} \Leftrightarrow \dfrac{1}{2^{n} } >\dfrac{1}{2^{n+1} }  \Leftrightarrow 2>1,\forall n\in \mathbb{N}^{*}$.\\
	Vậy $u_{n} =\dfrac{1}{2^{n} } ,\, \forall  n\in \mathbb{N}^{*}$ là dãy giảm.}
\end{ex}

% nguồn: Ok_TracNghiemLop11_Lan1
\begin{ex}%[1D2H1-4]%[ma:0165ab83]
	Trong các dãy số  $(u_{n})$ sau, dãy số nào là dãy số giảm?
	\choice 
		{$u_{n}=n^2$} 
		{$u_{n}=2n$} 
		{$u_{n}=n^3-1$} 
		 {\True $u_{n}=\dfrac{2n+1}{n-1}$}
	\loigiai{ 
		 $\forall n \in \mathbb{N}^*$, ta có:
		\begin{itemize}
			\item $n^2<(n+1)^2$ nên A  sai.
			\item $2n< 2(n+1)$ nên B  sai.
			\item $n^3-1<(n+1)^3-1$ nên C sai. 
		\end{itemize}
		Với $u_{n}=\dfrac{2n+1}{n-1}$  thì $u_{n+1}-u_{n}=\dfrac{-3}{n(n-1)}<0$ nên dãy $u_{n}=\dfrac{2n+1}{n-1}$ là dãy giảm.
	} 
\end{ex}

% nguồn: Ok_TracNghiemLop11_Lan1
\begin{ex}%[1D2H1-4]%[ma:4e7fc3be]
	Trong các dãy số  $(u_{n})$ sau, dãy số nào là dãy số tăng?
	\choice 
		{\True $u_{n} =\left(-1\right)^{2n}  \left(5^{n} +1\right),\, \forall n\in \mathbb{N}^{*} $}
		{$u_{n} =\dfrac{1}{\sqrt{n+1} +n}  ,\, \forall n\in \mathbb{N}^{*} $}
		{$u_{n} =\left(-1\right)^{n+1} \sin \dfrac{\pi }{n} ,\, \forall n\in \mathbb{N}^{*} $} 
		 {$u_{n} =\dfrac{n}{n^{2} +1} ,\, \forall n\in \mathbb{N}^{*} $}
	 \loigiai{
		Xét dãy số $u_{n} =\left(-1\right)^{2n} \left(5^{n} +1\right),\,  \forall n\in \mathbb{N}^{*}$.\\
		Có $u_{n} <u_{n+1} \Leftrightarrow 5^{n} +1<5^{n+1} +1\Leftrightarrow 1<5,\forall n\in  \mathbb{N}^{*}$.\\
	Vậy dãy số $(u_{n})$ là dãy  số tăng.}
\end{ex}

% nguồn: Ok_TracNghiemLop11_Lan1
\begin{ex}%[1D2H1-4]%[ma:18ad1215]
	Cho dãy  số $\left(u_{n}\right)$ biết $u_{n}=\dfrac{n^{2}+3}{2n^{2}-1}$ với $n\in N^{*}$. Tìm số hạng $u_{5}$.
	 \choice
		 {$u_{5}=\dfrac{7}{4}$}
		{$u_{5}=\dfrac{7}{9}$} 
		{$u_{5}=\dfrac{24}{51}$} 
		{\True $u_{5}=\dfrac{4}{7}$} 
	 \loigiai{
		$u_{5}=\dfrac{5^{2}+3}{2.5^{2}-1}=\dfrac{28}{49}=\dfrac{4}{7}$. 
	 }
\end{ex}

% nguồn: Ok_Toan11_TracNghiem_Lan2
\begin{ex}%[1D2V1-4]%[ma:53c84fe1]
	Dãy số nào sau đây  là dãy số giảm?
	\choice 
		{\True $u_n=\dfrac{5-3n}{2n+3}$} 
		 {$u_n=\dfrac{n-5}{4n+1}$}
		 {$u_n=2n^2+3$}
		 {$u_n=\cos (2n+1)$}
	\loigiai{ 
		\begin{itemize}
			\item Xét dãy $(u_n)$ với  $u_n=\dfrac{5-3n}{2n+3}$.\\
			Ta  có $u_{n+1}-u_n=\dfrac{5-3(n+1)}{2(n+1)+3}-\dfrac{5-3n}{2n+3}=\dfrac{-19}{(2n+3)(2n+5)} < 0, \forall n\in \mathbb{N}^{*}$.\\
			Vậy dãy số $(u_n)$ là  dãy số giảm.
			\item  Xét dãy $(u_n)$ với $u_n=\dfrac{n-5}{4n+1}$.\\
			Ta có $u_{n+1}-u_n=\dfrac{(n+1)-5}{4(n+1)+1}-\dfrac{n-5}{4n+1}=\dfrac{21}{(4n+5)(4n+1)}>0, \forall n\in  \mathbb{N}^{*}$.\\
			Vậy dãy số  $(u_n)$ là dãy số tăng.
			\item Xét dãy $(u_n)$ với  $u_n=2n^2+3$.\\
			Ta có $u_{n+1}-u_n=2(n+1)^2+3-(2n^2+3)=4n+1>0,  \forall n\in \mathbb{N}^{*}$.\\
			Vậy dãy số $(u_n)$ là dãy số  tăng.
			\item Xét dãy $(u_n)$ với $u_n=\cos (2n+1)$.\\ 
			Ta có $u_1\approx -0{,}98999$; $u_2\approx 0{,}28366$;  $u_4\approx -0{,}91113$.\\
			Vậy dãy số  $(u_n)$ không tăng, không giảm.
		\end{itemize}
	} 
\end{ex}

% nguồn: Ok_Toan11_TracNghiem_Lan2
\begin{ex}%[1D2V1-4]%[ma:450309ef]
	Trong các dãy  số $\left(u_n \right)$ cho bởi số hạng tổng quát $u_n$ sau, dãy số nào tăng?
	\choice 
		{$u_n=\dfrac{\sin  n}{n}$}
		{$u_n=\dfrac{\sqrt{n^2+1}}{2n+1}$} 
		{$u_n=\dfrac{3^n}{n^2}$} 
		{\True  $u_n=4n^3-3n^2+1$}
	 \loigiai{
		\begin{itemize}[*]
			\item Xét  dãy $\left(u_n \right)$ với $u_n=\dfrac{\sin n}{n}$.\\
			Với $n\in \left(k2\pi; \pi+k2\pi  \right), k\in \mathbb{N}\Rightarrow \sin n > 0\Rightarrow \dfrac{\sin n}{n} > 0$.\\
			Với $n\in \left(\pi+k2\pi; 2\pi+k2\pi  \right) $, $ k\in \mathbb{N}\Rightarrow \sin n < 0\Rightarrow \dfrac{\sin n}{n} < 0$. \\
			Suy ra dãy số $u_n=\dfrac{\sin n}{n}$ không tăng,  không giảm.
			\item Xét dãy $\left(u_n \right)$ với $u_n=\dfrac{\sqrt{n^2+1}}{2n+1}$.\\ 
			Ta có  $u_n=\dfrac{\sqrt{n^2+1}}{2n+1}=\sqrt{\dfrac{n^2+1}{\left(2n+1\right)^2}}$.\\
			Xét dãy $\left(v_n  \right)$ với $v_n=\dfrac{n^2+1}{\left(2n+1\right)^2}$.
			 \[v_{n+1}-v_n=\dfrac{n^2+2n+2}{4n^2+12n+9}-\dfrac{n^2+1}{4n^2+4n+1}=\dfrac{4n^2-2n-7}{\left(2n+3\right)^2\left(2n+1\right)^2}.
			\] 
			Do $v_{n+1}-v_n$  vừa nhận giá trị âm lẫn dương nên dãy số $\left(v_n \right)$ không tăng, không giảm $ \Rightarrow $ $ (u_n) $ không tăng, không giảm.
			\item Xét dãy $\left(u_n \right)$ với  $u_n=\dfrac{3^n}{n^2}$.\\
			$u_{n+1}-u_n=\dfrac{3.3^n}{\left(n+1\right)^2}-\dfrac{3^n}{n^2}=\dfrac{3^n \left(2n^2-2n-1\right)}{\left(n+1\right)^2n^2}$.\\ 
			Do $u_{n+1}-u_n$ nhận giá trị âm lẫn dương  nên dãy đã cho không tăng, không giảm.
			\item Xét dãy  $\left(u_n \right)$ với $u_n=4n^3-3n^2+1$.\\
			$u_{n+1}-u_n= 4(n+1)^3-3(n+1)^2+1-(4n^3-3n^2+1)=4(n+1)^3-4n^3 -3(n+1)^2 + 3n^2 >0  $.\\
			{\centering (Do $ n\in \mathbb{N}^{*} $ nên $ 4(n+1)^3-4n^3 >0 $ và $  -3(n+1)^2 + 3n^2 >0 $)}\\
			Do đó $u_n=4n^3-3n^2+1$ là dãy số tăng. 
		\end{itemize}
	} 
\end{ex}

% nguồn: Ok_Toan11_TracNghiem_Lan2
\begin{ex}%[1D2V1-4]%[ma:eb56bdeb]
	Cho dãy số $(u_n)$ biết $\heva{&u_1=1\\ 
	&u_{n+1}=au_n+1}$. Tìm  tất cả các giá trị của $a$ để $(u_n)$ tăng.
	\choice 
		 {$a < 0$}
		{$a\leq  0$}
		{\True $a  > 0$}
		{$a >  1$}
	\loigiai{ 
		 Ta có $u_2=a\cdot 1 + 1= a+1$.
		Xét hiệu  $u_{n+1}-u_n=a\left(u_n-u_{n-1}\right)$.\\
		Tương  tự, ta có
		\begin{itemize}
			\item $u_n-u_{n-1}=a(u_{n-1}-u_{n_2})$; 
			\item  $\ldots$;
			 \item $u_3-u_2=a(u_2-u_1)$;
			\item  $u_2-u_1=a$.
		\end{itemize}
		 Từ đó, ta được $u_{n+1}-u_{n}=a^{n-1}$.\\
	Để dãy số $(u_n)$ là dãy số tăng thì $a^{n-1}>0$, $\forall  n\in\mathbb{N}^*$, tương đương $a>0$.}
\end{ex}

% nguồn: Ok_Toan11_TracNghiem_Lan2
\begin{ex}%[1D2V1-4]%[ma:b49c6251]
	Cho dãy số($u_n$) biết $u_n=\dfrac{3^n}{an}$. Tìm tất cả các giá trị của $a$ để dãy số tăng. 
	 \choice
		{$\forall  a < 0$}
		{Không tồn tại  $a$}
		{$\forall a \in \mathbb{R}^*$} 
		{\True $a >  0$}
	\loigiai{Ta có  $u_{n+1}-u_n = \dfrac{3^{n+1}}{an+a} - \dfrac{3^n}{an} = \dfrac{3^n\cdot a\cdot(2n-1) }{a^2(n^2+n) }$, $\forall n \in \mathbb{N}^*$.\\
	Để  dãy số tăng thì $u_{n+1}-u_n = \dfrac{a\cdot3^n(2n-1) }{a^2(n^2+n) } > 0$, $\forall n \in \mathbb{N}^* \Leftrightarrow a > 0$.}
\end{ex}

% nguồn: Ok_Toan11_TracNghiem_Lan2
\begin{ex}%[1D2V1-4]%[ma:10991575]
	Trong các dãy số $(u_n)$ cho bởi số hạng tổng quát $u_n$ sau, dãy số nào là dãy số tăng? 
	\choice 
		 {$u_n=\dfrac{2}{3^n}$}
		 {$u_n=\dfrac{3}{n}$}
		{\True  $u_n=2^n$}
		{$u_n=\left(-2\right)^n$} 
	\loigiai{ 
		\begin{itemize}
			\item Xét  dãy $(u_n)$ với $u_n=\dfrac{2}{3^n}$.\\
			Ta có $\dfrac{u_{n+1}}{u_n}=\dfrac{2}{3^{n+1}}:\dfrac{2}{3^n}=\dfrac{1}{3}<1, \forall n\in  \mathbb{N}^{*}$.\\
			Vậy dãy số $(u_n)$ là  dãy số giảm.
			\item Xét dãy $(u_n)$  với $u_n=\dfrac{3}{n}$.\\
			Ta có $u_{n+1}-u_n=\dfrac{3}{n+1}-\dfrac{3}{n}=-\dfrac{3}{n(n+1)}<0, \forall n\in \mathbb{N}^{*}$.\\ 
			Vậy  dãy số $(u_n)$ là dãy số giảm.
			\item Xét dãy $(u_n)$ với $u_n=2^n$.\\ 
			Ta có $\dfrac{u_{n+1}}{u_n}=\dfrac{2^{n+1}}{2^n}=2>1, \forall n\in  \mathbb{N}^{*}$.\\
			Vậy dãy số $(u_n)$ là dãy số tăng. 
			\item Xét  dãy $(u_n)$ với $u_n=\left(-2\right)^n$.\\
			Ta có $u_1=-2$;  $u_2=4$; $u_3=-8$.\\
			 Vậy dãy số $(u_n)$ là dãy số không tăng, không giảm.
		\end{itemize}
	} 
\end{ex}

% nguồn: Ok_Toan11_TracNghiem_Lan2
\begin{ex}%[1D2V1-4]%[ma:e251eb85]
	Cho dãy số  $(u_n)$ biết $u_n=\sqrt{3n+2}-\sqrt{3n+1}$. Mệnh đề nào sau đây đúng?
	 \choice
		{Dãy  số tăng}
		{\True Dãy số giảm} 
		{Dãy  số không tăng, không giảm}
		{Dãy  số vừa tăng vừa giảm}
	 \loigiai{
		Ta có  \[u_n=\sqrt{3n+2}-\sqrt{3n+1}=\dfrac{1}{\sqrt{3n+2}+\sqrt{3n+1}}.\]
		\[3n+5>3n+2>0  \Rightarrow \sqrt{3n+5}>\sqrt{3n+2}, \forall n\in \mathbb{N}^{*}.\]
		\[3n+4>3n+1>0 \Rightarrow  \sqrt{3n+4}>\sqrt{3n+1}, \forall n\in \mathbb{N}^{*}.\]
		\begin{align*}
			&\Rightarrow \sqrt{3n+5} +\sqrt{3n+4}>  \sqrt{3n+2}+\sqrt{3n+1}\\
			&\Rightarrow \dfrac{1}{\sqrt{3n+5}  +\sqrt{3n+4}}< \dfrac{1}{\sqrt{3n+2}+\sqrt{3n+1}}\\
			&\Rightarrow \dfrac{1}{\sqrt{3n+5} +\sqrt{3n+4}}- \dfrac{1}{\sqrt{3n+2}+\sqrt{3n+1}}<0,  \forall n\in \mathbb{N}^{*}.\\
			&\Rightarrow u_{n+1}-u_n <0, \forall n\in \mathbb{N}^{*}. 
		\end{align*}
		Do đó dãy $u_n=\sqrt{3n+2}-\sqrt{3n+1}$ là số  giảm.
	} 
\end{ex}

% nguồn: Ok_Toan11_TracNghiem_Lan2
\begin{ex}%[1D2V1-4]%[ma:6f0b0cc5]
	Cho dãy  số $(u_n)$ biết $u_n=n-\sqrt{n^2+1}$. Mệnh đề nào sau đây đúng?
	\choice 
		{\True Dãy số tăng} 
		{Dãy số  giảm}
		 {Dãy số không tăng, không giảm}
		{Các số hạng đều dương} 
	\loigiai{ 
		Ta  có \[u_n=n-\sqrt{n^2+1}=\dfrac{-1}{n+\sqrt{n^2+1}}.\]
		\[n+1>n>0,  \forall n\in \mathbb{N}^{*}.\]
		\[(n+1)^2+1>n^2+1>0 \Rightarrow \sqrt{(n+1)^2+1}>\sqrt{n^2+1}, \forall n\in \mathbb{N}^{*}.\] 
		 \allowdisplaybreaks
		\begin{align*}
			 &\Rightarrow n+1 +\sqrt{(n+1)^2+1}> n +\sqrt{n^2+1}\\
			 &\Rightarrow \dfrac{1}{\sqrt{ n+1 +\sqrt{(n+1)^2+1}}} < \dfrac{1}{n +\sqrt{n^2+1}}\\
			&\Rightarrow \dfrac{1}{\sqrt{  n+1 +\sqrt{(n+1)^2+1}} }- \dfrac{1}{n +\sqrt{n^2+1}}<0, \forall n\in \mathbb{N}^{*}.\\
			&\Rightarrow \dfrac{-1}{\sqrt{ n+1 +\sqrt{(n+1)^2+1}} }- \dfrac{-1}{n +\sqrt{n^2+1}}>0, \forall  n\in \mathbb{N}^{*}.\\
			&\Rightarrow u_{n+1}-u_n >0, \forall  n\in \mathbb{N}^{*}.
		\end{align*}
		Do đó dãy $u_n=n-\sqrt{n^2+1}$  là số tăng.
	} 
\end{ex}

% nguồn: Ok_Toan11_TracNghiem_Lan2
\begin{ex}%[1D2V1-4]%[ma:e853672f]
	Cho dãy số $\left(u_n \right)$ biết $\heva{&u_1=3  \\ &u_{n+1}=\dfrac{3u_n}{3+u_n}.}$ Mệnh đề nào sau đây đúng?
	 \choice
		{Dãy số  tăng}
		{\True Dãy số giảm} 
		{Dãy số không  tăng, không giảm}
		{Có  $u_{10}=2$}
	\loigiai{Ta có $u_1 > u_2 > u_3$. Dự đoán dãy số đã cho giảm, ta  chứng minh bằng quy nạp.\\
		Từ giả thiết thì $u_n  > 0$, $\forall n \in \mathbb{N}^*$.\\
		Giả sử $u_k < u_{k-1}$, $k \ge 2$. Ta  chứng minh $u_{k+1} < u_k$. \\
	Thật vậy $$u_{k+1}-u_k = \dfrac{3u_k}{3+u_k} - \dfrac{3u_{k-1}}{3+u_{k-1}} = \dfrac{9(u_k-u_{k-1}) }{(3+u_k) (3+u_{k-1}) } < 0 \Leftrightarrow u_{k+1} < u_k.$$ Vậy dãy đã  cho là dãy giảm.}
\end{ex}

% nguồn: Ok_Toan11_TracNghiem_Lan2
\begin{ex}%[1D2V1-4]%[ma:a8b45390]
	Trong các dãy số dưới đây, dãy số nào  là dãy số giảm?
	 \choice
		{$u_n=n^2$} 
		 {\True $u_n=\dfrac{1}{n}-3$}
		 {$u_n=3n$}
		 {$u_n=n^3-2$}
	\loigiai{ 
		\begin{itemize}
			 \item Xét dãy $(u_n)$ với $u_n=n^2$.\\
			Ta có $u_{n+1}-u_n=(n+1)^2-n^2=2n+1>  0,\forall n\in \mathbb{N}^{*}$.\\
			Vậy dãy số $(u_n)$ là dãy số tăng. 
			\item Xét  dãy $(u_n)$ với $u_n=\dfrac{1}{n}-3$.\\
			Ta có $u_{n+1}-u_n=\dfrac{1}{n+1}-\dfrac{1}{n}=\dfrac{-1}{n(n+1)} < 0,\forall n\in  \mathbb{N}^{*}$.\\
			Vậy dãy số $(u_n)$ là dãy số giảm. 
			\item Xét dãy $(u_n)$  với $u_n=3n$.\\
			Ta có $u_{n+1}-u_n=3(n+1)-3n=3>  0,\forall n\in \mathbb{N}^{*}$.\\
			Vậy dãy số $(u_n)$  là dãy số tăng.
			 \item Xét dãy $(u_n)$ với $u_n=n^3-2$.\\
			 Ta có $u_{n+1}-u_n=(n+1)^3-n^3> 0,\forall n\in \mathbb{N}^{*}$.\\
			Vậy dãy số $(u_n)$ là dãy số  tăng.
		\end{itemize}
	 }
\end{ex}

% nguồn: Ok_Toan11_TracNghiem_Lan2
\begin{ex}%[1D2V1-4]%[ma:d41a607b]
	Trong các dãy số sau, dãy  số nào là dãy số giảm?
	\choice 
		{$u_n=\dfrac{n-3}{n+1}$} 
		 {$u_n=\dfrac{n}{2}$}
		{\True $u_n=\dfrac{2}{n^2}$} 
		 {$u_n=\dfrac{\left(-1\right)^n}{3^n}$}
	 \loigiai{
		\begin{itemize}
			\item Xét dãy $(u_n)$ với $u_n=\dfrac{n-3}{n+1}$.\\ 
			Ta có $u_{n+1}-u_n=\dfrac{n-2}{n+2}-\dfrac{n-3}{n+1}=\dfrac{4}{\left(n+1\right)\left(n+2\right)} > 0,  \forall n\in \mathbb{N}^*$.\\
			Vậy dãy số $(u_n)$  là dãy số tăng.
			\item Xét dãy $(u_n)$ với $u_n=\dfrac{n}{2}$.\\ 
			Ta có $u_{n+1}-u_n=\dfrac{n+1}{2}-\dfrac{n}{2}=\dfrac{1}{2} > 0, \forall  n\in \mathbb{N}^*$.\\
			Vậy dãy số $(u_n)$  là dãy số tăng.
			\item Xét dãy  $(u_n)$ với $u_n=\dfrac{2}{n^2}$.\\
			 Ta có $\dfrac{u_{n+1}}{u_n}=\dfrac{n^2}{\left(n+1\right)^2} < \dfrac{n^2}{n^2}=1,\forall n\in \mathbb{N}^{*}$.\\
			Vậy dãy số $(u_n)$ là  dãy số giảm.
			\item Xét  dãy $(u_n)$ với $u_n=\dfrac{\left(-1\right)^n}{3^n}$.\\
			 Ta có $u_1=\dfrac{-1}{3};$ $u_2=\dfrac{1}{9};$ $u_3=\dfrac{-1}{27}$.\\
			Vậy  dãy số $(u_n)$ là dãy số không tăng, không giảm.
		\end{itemize}
	} 
\end{ex}

% nguồn: Ok_Toan11_TracNghiem_Lan2
\begin{ex}%[1D2V1-4]%[ma:a57d3a0f]
	Trong các dãy số sau, dãy số nào là dãy  số giảm?
	 \choice
		{\True  $u_n=\dfrac{3}{n^2}$}
		{$u_n=\dfrac{n-3}{n+1}$} 
		 {$u_n=\dfrac{n}{2}$}
		 {$u_n=\dfrac{\left(-1\right)^n}{3^n}$}
	\loigiai{ 
		\begin{itemize}
			 \item Xét dãy $(u_n)$ với $u_n=\dfrac{3}{n^2}$.\\
			Ta có $\dfrac{u_{n+1}}{u_n}=\dfrac{n^2}{(n+1)^2} <  \dfrac{n^2}{n^2}=1,\forall n\in \mathbb{N}^{*}$.\\
			Vậy dãy  số $(u_n)$ là dãy số giảm.
			\item Xét dãy  $(u_n)$ với $u_n=\dfrac{n-3}{n+1}$.\\
			Ta có $u_{n+1}-u_n=\dfrac{n-2}{n+2}-\dfrac{n-3}{n+1}=\dfrac{4}{(n+1)(n+2)} >  0$ $\forall n\in \mathbb{N}^*$.\\
			Vậy dãy  số $(u_n)$ là dãy số tăng.
			\item Xét dãy $(u_n)$ với  $u_n=\dfrac{n}{2}$.\\
			Ta có $u_{n+1}-u_n=\dfrac{n+1}{2}-\dfrac{n}{2}=\dfrac{1}{2} >  0$ $\forall n\in \mathbb{N}^*$.\\
			Vậy  dãy số $(u_n)$ là dãy số tăng.
			\item Xét dãy $(u_n)$ với $u_n=\dfrac{\left(-1\right)^n}{3^n}$.\\ 
			Ta có $u_1=\dfrac{-1}{3};$  $u_2=\dfrac{1}{9};$ $u_3=\dfrac{-1}{27}$.\\
			Vậy dãy số $(u_n)$ là dãy số không tăng, không giảm. 
		\end{itemize}
	} 
\end{ex}

% nguồn: Ok_Toan11_TracNghiem_Lan2
\begin{ex}%[1D2V1-4]%[ma:01675698]
	Cho dãy số $(u_n)$ biết $u_n=\dfrac{an+2}{3n+1}$. Tìm tất cả các giá trị của $ a $ để dãy số  tăng.
	\choice 
		{$a=6$} 
		{\True $a >  6$}
		 {$a < 6$}
		 {$a\ge 6$}
	\loigiai{ 
		 Ta có $u_{n+1}-u_n=\dfrac{an+a+2}{3n+4}-\dfrac{an+2}{3n+1}=\dfrac{a-6}{\left(3n+4\right)\left(3n+1\right)},$, $ \forall n\in \mathbb{N}^{*}$.\\
		Để dãy số tăng thì $u_{n+1}-u_n=\dfrac{a-6}{\left(3n+4\right)\left(3n+1\right)}  > 0, \forall n\in \mathbb{N}^{*} \Leftrightarrow a > 6$.
	 }
\end{ex}

% nguồn: Ok_Toan11_TracNghiem_Lan2
\begin{ex}%[1D2V1-4]%[ma:0f4e485a]
	Cho dãy số $(u_n)$ biết $u_n=2^n-an$. Tìm tất cả các giá trị của $a$ để dãy số tăng. 
	\choice 
		{$a=2$} 
		{$a > 2$} 
		 {\True $a < 2$}
		{$a \geq 2$} 
	 \loigiai{
		Ta có $u_{n+1}-u_n=2^{n+1}-an-a-2^n+a n=2^n-a, \forall n  \in \mathbb{N}^*$.\\
		Dãy số tăng thì $u_{n+1}-u_n=2^n-a > 0, \forall n \in \mathbb{N}^* \Leftrightarrow a < 2^n, \forall  n \in \mathbb{N}^* \Leftrightarrow a < 2, \forall n \in \mathbb{N}^*$.
	 }
\end{ex}

% nguồn: Ok_Toan11_TracNghiem_Lan2
\begin{ex}%[1D2V1-4]%[ma:af5dda05]
	Cho dãy số $\left(u_n \right)$ biết  $u_n=\dfrac{1}{n+1}+\dfrac{1}{n+2}+\cdots+\dfrac{1}{n+n}$. Mệnh đề nào sau đây đúng?
	 \choice
		 {\True Dãy số tăng}
		{Dãy số giảm} 
		{Dãy số không tăng,  không giảm}
		 {Có hữu hạn số hạng}
	\loigiai{Xét hiệu $u_{n+1}-u_n = \dfrac{1}{2n+1} + \dfrac{1}{2n+2} - \dfrac{1}{n+1}  = \dfrac{4n^2+3n+1}{2(2n+1) (n+1)^2} > 0$, $\forall n \in \mathbb{N}^*$.}
\end{ex}

% nguồn: Ok_Toan11_TracNghiem_Lan2
\begin{ex}%[1D2V1-4]%[ma:36e61491]
	Cho dãy số $(u_n)$ biết $u_n=\dfrac{3^n}{an}$. Tìm tất cả các giá trị của $ a $ để dãy số tăng. 
	 \choice
		{$  a < 0$}
		{Không tồn  tại a}
		{$a\in  \mathbb{R}^{*}$}
		 {\True $a > 0$}
	\loigiai{ 
		Ta có $u_{n+1}-u_n=\dfrac{3^{n+1}}{an+a}-\dfrac{3^n}{an}=\dfrac{a.3^n \left(2n-1\right)}{a^2\left(n^2+n\right)}, \forall n\in  \mathbb{N}^{*}$.\\
		Để dãy số tăng thì $u_{n+1}-u_n=\dfrac{a.3^n \left(2n-1\right)}{a^2\left(n^2+n\right)} > 0$,  $ \forall n\in \mathbb{N}^{*} \Leftrightarrow a > 0$.
	 }
\end{ex}

% nguồn: Ok_Toan11_TracNghiem_Lan2
\begin{ex}%[1D2V1-4]%[ma:075f300d]
	Trong các dãy số $(u_n)$ cho  bởi số hạng tổng quát $u_n$ sau, dãy số nào là dãy số tăng?
	\choice 
		 {$u_n=\dfrac{1}{2^n}$}
		 {\True $u_n=\dfrac{3n-1}{n+1}$}
		{$u_n=1-n^2$} 
		{$u_n=\dfrac{1}{\sqrt{n+2}}$} 
	\loigiai{ 
		\begin{itemize}
			\item Xét  dãy $(u_n)$ với $u_n=\dfrac{1}{2^n}$.\\
			Ta có $\dfrac{u_{n+1}}{u_n}=\dfrac{1}{2^{n+1}}:\dfrac{1}{2^n}=\dfrac{1}{2}<1, \forall  n\in \mathbb{N}^{*}$.\\
			Vậy dãy số $(u_n)$ là dãy số giảm. 
			 \item Xét dãy $(u_n)$ với $u_n=\dfrac{3n-1}{n+1}$.\\
			Ta có $u_{n+1}-u_n=\dfrac{3(n+1)-1}{(n+1)+1}-\dfrac{3n-1}{n+1}=\dfrac{4}{(n+2)(n+1)}>0, \forall n\in  \mathbb{N}^{*}$.\\
			Vậy dãy số $(u_n)$ là dãy  số tăng.
			\item  Xét dãy $(u_n)$ với $u_n=1-n^2$.\\
			Ta có $u_{n+1}-u_n=1-(n+1)^2-(1-n^2)=-2n-1<0, \forall n\in \mathbb{N}^{*}$.\\ 
			Vậy dãy  số $(u_n)$ là dãy số giảm.
			\item Xét dãy $(u_n)$ với $u_n=\dfrac{1}{\sqrt{n+2}}$.\\ 
			Ta có $\dfrac{u_{n+1}}{u_n}=\dfrac{1}{\sqrt{n+3}}:\dfrac{1}{\sqrt{n+2}}=\sqrt{\dfrac{n+2}{n+3}}=\sqrt{1-\dfrac{1}{n+3}}<1, \forall n\in  \mathbb{N}^{*}$.\\
			Vậy dãy số  $(u_n)$ là dãy số giảm.
		\end{itemize}
	} 
\end{ex}

% nguồn: Ok_Toan11_TracNghiem_Lan2
\begin{ex}%[1D2V1-4]%[ma:0547e2ae]
	Cho dãy số $(u_n)$ biết $u_n=\dfrac{1}{n+1}+\dfrac{1}{n+2}+\cdots+\dfrac{1}{n+n}$. Mệnh đề nào sau đây  đúng?
	\choice 
		{\True Dãy số tăng} 
		{Dãy số giảm} 
		{Dãy số không tăng, không giảm} 
		{Có  hữu hạn số hạng}
	\loigiai{ 
		 Ta có
		\allowdisplaybreaks 
		\begin{eqnarray*}
			u_{n+1}&=&\dfrac{1}{n+1+1}+\dfrac{1}{n+1+2}+\ldots+\dfrac{1}{n+1+n+1} \\ 
			&=&  \dfrac{1}{n+2}+\dfrac{1}{n+3}+\ldots+\dfrac{1}{n+n+2}.
		\end{eqnarray*}
		 Khi đó
		\allowdisplaybreaks 
		\begin{eqnarray*}
			u_{n+1}-u_n&=&-\dfrac{1}{n+1}+\dfrac{1}{2n+1}+\dfrac{1}{2n+2}  \\
			&=& \dfrac{1}{2(n+1)(2n+1)}  >0, \forall n\in\mathbb{N}^*.
		\end{eqnarray*}
		 Suy ra $u_{n+1}>u_n$.\\
		Vậy dãy số $(u_n)$ là dãy số  tăng
	} 
\end{ex}

% nguồn: Ok_Toan11_TracNghiem_Lan2
\begin{ex}%[1D2V1-4]%[ma:8015f930]
	Trong  các dãy số $(u_n)$ cho bởi số hạng tổng quát $u_n$ sau, dãy số nào là dãy số giảm?
	\choice 
		{\True $u_n=\dfrac{1}{2^n}$} 
		{$u_n=\dfrac{3n-1}{n+1}$} 
		{$u_n=n^2$} 
		{$u_n=\sqrt{n+2}$} 
	\loigiai{ 
		\begin{itemize}
			\item Xét dãy $(u_n)$ với  $u_n=\dfrac{1}{2^n}$.\\
			Ta có $\dfrac{u_{n+1}}{u_n}=\dfrac{1}{2}<1, \forall  n\in \mathbb{N}^{*}$.\\
			 Vậy dãy số $(u_n)$ là dãy số giảm.
			\item Xét  dãy $(u_n)$ với $u_n=\dfrac{3n-1}{n+1}$.\\
			Ta có $u_{n+1}-u_n=\dfrac{3(n+1)-1}{(n+1)+1}-\dfrac{3n-1}{n+1}=\dfrac{4}{(n+2)(n+1)}>0, \forall n\in  \mathbb{N}^{*}$.\\
			Vậy dãy$(u_n)$ là dãy số  tăng.
			\item Xét dãy $(u_n)$  với $u_n=n^2$.\\
			 Ta có $u_{n+1}-u_n=(n+1)^2-n^2=2n+1>0, \forall n\in \mathbb{N}^{*}$.\\
			Vậy dãy số $(u_n)$ là dãy số tăng. 
			 \item Xét dãy $(u_n)$ với $u_n=\sqrt{n+2}$.\\
			Ta  có $\dfrac{u_{n+1}}{u_n}=\dfrac{\sqrt{n+3}}{\sqrt{n+2}}=\sqrt{1+\dfrac{1}{n+2}}>\sqrt{1}=1, \forall n\in \mathbb{N}^{*}$.\\
			Vậy  dãy số $(u_n)$ là dãy số tăng.
		\end{itemize}
	} 
\end{ex}

% nguồn: Ok_TracNghiemLop11_Lan1
\begin{ex}%[1D2V1-4]%[ma:938e873c]
	Cho  dãy số$\left(u_n\right)$ xác định bởi:
	$\heva{&u_{1}=2;~u_{2}=3  \\
	&u_{n+2}=3u_{n+1}-2u_{n} \text{ với } n\ge 1\\}$ 
	 Khẳng định nào sau đây sai?
	 \choice
		 {${{u}_{n}}={{2}^{n-1}}+1$}
		{$\left( {{u}_{n}} \right)$là dãy số tăng} 
		{Năm số hạng đầu của  dãy số là: $2$, $3$, $5$, $9$, $17$}
		{\True  $u_{n}=\dfrac{n^{2}+5}{3}$}
	 \loigiai{
		Từ công thức truy  hồi ta có: $u_{1}=2$; $u_{2}=3$; $u_{3}=5$; $u_{4}=9$; $u_{5}=17$,... \\
		 Xét: $u_{n+2}=3u_{n+1}-2u_{n}$
		$\Rightarrow u_{n}=\dfrac{3u_{n+1}-u_{n+2}}{2}$.  \\
		Vậy: $u_{1}=2=1+1=2^{0}+1$. \\ 
		 $u_{2}=3=2+1=2^{1}+1$. \\
		$u_{3}=5=4+1=2^{2}+1$.  \\
		$u_{4}=9=8+1=2^{3}+1$.  \\
		$u_{5}=17=16+1=2^{4}+1$  \\
		 ..... \\
	 Vậy công thức tổng quát của dãy số là: $u_{n}=2^{n-1}+1$.}
\end{ex}

% nguồn: Ok_Toan11_TracNghiem_Lan2
\begin{ex}%[1D2C1-4]%[ma:7fdf1e4a]
	Cho dãy số  $(u_n)$ biết $\heva{& u_1=3 \\
	 & u_{n+1}=\dfrac{3u_n}{3+u_n}}$. Mệnh đề nào sau đây đúng?
	\choice 
		 {Dãy số tăng}
		 {\True Dãy số giảm}
		{Dãy số  không tăng, không giảm}
		{Có $u_{10}=2$} 
	\loigiai{ 
		Ta có $u_1 > u_2 > u_3$. Dự đoán dãy số đã cho giảm, ta chứng minh bằng quy nạp.\\ 
		Từ giả thiết thì $u_n >  0, \forall n\in \mathbb{N}^{*}$.\\
		Giả sử $u_k <  u_{k-1}, k\ge 2$. Ta chứng minh $u_{k+1} < u_k$.\\
		Thật vậy $u_{k+1}-u_k=\dfrac{3u_k}{3+u_k}-\dfrac{3u_{k-1}}{3+u_{k-1}}=\dfrac{9\left(u_k-u_{k-1} \right)}{\left(3+u_k \right)\left(3+u_{k-1} \right)} < 0$.\\ 
		Suy  ra $u_{k+1} < u_k$.\\
		 Vậy dãy số $(u_n)$ là dãy số giảm.
	 }
\end{ex}

% nguồn: Ok_Toan11_TracNghiem_Lan2
\begin{ex}%[1D2C1-4]%[ma:11e0c927]
	Cho  dãy số $(u_n)$ biết $\heva{& u_1=1 \\
	& u_{n+1}=\sqrt{u_n^2+3}}$. Mệnh đề nào sau đây đúng? 
	\choice 
		{\True Dãy  số tăng}
		 {Dãy số giảm}
		{Dãy số không  tăng, không giảm}
		{Dãy  số có $u_5=13$}
	 \loigiai{
		Ta có $0< u_1 < u_2 < u_3$.  Dự đoán dãy số đã cho tăng, ta chứng minh bằng quy nạp.\\
		Từ giả thiết thì $u_n > 0$, $\forall  n\in \mathbb{N}^{*}$.\\
		Giả sử $u_k > u_{k-1}, k\ge 2$. Ta  chứng minh $u_{k+1} > u_k$.\\
		Thật vậy $u_{k+1}-u_k=\sqrt{u_k^2+3}-\sqrt{u_{k-1}^2+3}=\dfrac{\left(u_k-u_{k-1} \right)\left(u_k+u_{k-1} \right)}{\sqrt{u_k^2+3}+\sqrt{u_{k-1}^2+3}}  > 0$.\\
		Suy ra $u_{k+1} > u_k$.\\ 
		Vậy  dãy số $(u_n)$ cho là dãy số tăng.
	} 
\end{ex}


%%%%%%%%%%%%%%%% ĐÚNG/SAI %%%%%%%%%%%%%%%%

% nguồn: Ok_dungsai_lop11_Lan1
\begin{ex}%[1D2H1-4]%[ma:f04da51b]
 Các mệnh  đề sau đúng hay sai?
 \choiceTF 
  {Dãy số $u_n=\pi$  tăng}
  {\True Dãy số $u_n=\dfrac{2023}{n}$ không tăng} 
   {\True Dãy số $u_n=n^3$ tăng}
  {\True Dãy số  $u_n=n!$ có $u_{n+1}=(n+1)u_n$}
 \loigiai{ 
  \begin{itemchoice}
   \itemch Mệnh đề \lq\lq Dãy số $u_n=\pi$ tăng\rq\rq\, sai. 
   \itemch  Mệnh đề \lq\lq Dãy số $u_n=\dfrac{2023}{n}$ không tăng\rq\rq\, đúng.
   \itemch Mệnh đề \lq\lq Dãy số $u_n=n^3$ tăng\rq\rq\,  đúng.
   \itemch Mệnh đề \lq\lq Dãy  số $u_n=n!$ có $u_{n+1}=(n+1)u_n$\rq\rq\, đúng.
  \end{itemchoice}
  }
\end{ex}

% nguồn: Ok_dungsai_lop11_Lan1
\begin{ex}%[1D2H1-4]%[ma:274a8c26]
 Cho dãy số $(u_n)$ có số hạng tổng quát $u_n=1-\dfrac{1}{n}$. 
 \choiceTF 
   {\True $u_3=\dfrac{2}{3}$}
   {\True $u_7-u_8=\dfrac{1}{56}$}
  {$u_{n+1}-u_n=-\dfrac{1}{n(n+1)}$} 
   {\True Dãy số $(u_n)$ là dãy số tăng}
  \loigiai{
  \begin{itemchoice}
    \itemch Đúng vì $u_3=1-\dfrac{1}{3}=\dfrac{2}{3}$.
   \itemch Sai  vì $u_7-u_8=-\dfrac{1}{56}$
   \itemch Sai vì $u_{n+1}-u_n=\left(1-\dfrac{1}{n+1}\right)-\left(1-\dfrac{1}{n}\right)=\dfrac{1}{n}-\dfrac{1}{n+1}=\dfrac{1}{n(n+1)}>0, \forall n \geq  1$.
   \itemch Đúng vì $u_{n+1}-u_n=\left(1-\dfrac{1}{n+1}\right)-\left(1-\dfrac{1}{n}\right)=\dfrac{1}{n}-\dfrac{1}{n+1}=\dfrac{1}{n(n+1)}>0,  \forall n \geq 1$.\\
   Suy ra $u_{n+1}>u_n,  \forall n \in \mathbb{N}^*$.\\
   Vậy dãy số $(u_n)$ là dãy  số tăng.
  \end{itemchoice}
  }
\end{ex}

% nguồn: Ok_dungsai_lop11_Lan1
\begin{ex}%[1D2H1-4]%[ma:d587cf49]
 Cho dãy số $(u_n)$ có số hạng tổng quát $u_n=1-\dfrac{1}{n}$.  Khi đó
  \choiceTF
  {\True $u_3=\dfrac{2}{3}$} 
   {$u_7-u_8=\dfrac{1}{56}$}
   {$u_{n+1}-u_n=-\dfrac{1}{n(n+1)}$}
  {\True Dãy số $(u_n)$  là dãy số tăng}
 \loigiai{\begin{itemchoice}
  \itemch \textbf{Đúng.} Ta có $u_3=\dfrac{2}{3}$. 
  \itemch \textbf{Sai.} Ta  có $u_7-u_8=\left(1-\dfrac{1}{7}\right)-\left(1-\dfrac{1}{8}\right)=-\dfrac{1}{56}$.
  \itemch \textbf{Sai.}Ta có $u_{n+1}-u_n=\left(1-\dfrac{1}{n+1}\right)-\left(1-\dfrac{1}{n}\right)=\dfrac{1}{n}-\dfrac{1}{n+1}=\dfrac{1}{n(n+1)}$. 
  \itemch \textbf{Đúng.} Do $u_{n+1}-u_n=\dfrac{1}{n(n+1)}>0$, $\forall n\in\mathbb{N}^*$ nên $u_{n+1}>u_n$, $\forall n\in\mathbb{N}^*$, suy ra $(u_n)$  là dãy số tăng.
\end{itemchoice}}
\end{ex}

% nguồn: Ok_dungsai_lop11_Lan1
\begin{ex}%[1D2H1-4]%[ma:721dc9df]
 Các mệnh đề sau đúng hay sai? 
  \choiceTF
  {\True Dãy số $u_n = (-1)^{2n +  1}\cdot 3^n$ là dãy giảm}
  {Dãy số $u_n = \dfrac{1}{n} - \dfrac{1}{n + 1}$ là  dãy không tăng}
  {\True Dãy số $u_n = 3n^2  - n^3$ là dãy không giảm}
  {Dãy số $u_n = \sqrt{n + 1} - \sqrt{n}$  là dãy tăng}
  \loigiai{
  \begin{itemchoice}
   \itemch $u_n = (-1)^{2n + 1}\cdot 3^n = -3^n$ là dãy  giảm suy ra mệnh đề \lq\lq Dãy số $u_n = (-1)^{2n + 1}\cdot 3^n$ là dãy giảm\rq\rq\, đúng.
   \itemch $u_n = \dfrac{1}{n} - \dfrac{1}{n + 1} = \dfrac{1}{n(n + 1)}$ là dãy giảm suy ra mệnh đề \lq\lq Dãy số $u_n = \dfrac{1}{n} - \dfrac{1}{n + 1}$  là dãy không tăng\rq\rq\, sai.
   \itemch $u_n = 3n^2 - n^3$ không là dãy đơn điệu suy ra mệnh đề \lq\lq Dãy số $u_n =  3n^2 - n^3$ là dãy không giảm\rq\rq\, đúng.
   \itemch $u_n = \sqrt{n + 1} - \sqrt{n} = \dfrac{1}{\sqrt{n+1} + \sqrt{n}}$ là dãy giảm suy ra mệnh đề \lq\lq Dãy số $u_n = \sqrt{n + 1} - \sqrt{n}$ là dãy  tăng\rq\rq\, sai.
  \end{itemchoice}
 } 
\end{ex}

% nguồn: Ok_dungsai_lop11_Lan1
\begin{ex}%[1D2H1-4]%[ma:ae3fee6e]
 Trong các dãy số $(u_n)$ cho bởi số hạng tổng quát $u_n$ sau, dãy số nào là  dãy số tăng?
  \choiceTF
  {Dãy  số $u_n=\dfrac{1}{2^n}$ tăng}
  {\True Dãy số  $u_n=\dfrac{1}{n}$ không tăng}
  {\True Dãy số  $u_n=\dfrac{2n-1}{n+1}$ giảm}
   {Dãy số $u_n=\dfrac{n+5}{3n+1}$ tăng}
 \loigiai{ 
 \begin{itemchoice}
   \itemch Xét dãy số $(u_n)$ có $u_n=\dfrac{1}{2^n}$. Với mọi $n\in \mathbb{N}^*$, ta có $\dfrac{u_{n+1}}{u_n} = \dfrac{2^n}{2^{n+1}}=\dfrac{1}{2}<1$, tức là $u_{n+1}<u_n$ cho nên dãy số $(u_n)$ đã cho là dãy số giảm.\\
  Mệnh đề \lq\lq Dãy số $u_n=\dfrac{1}{2^n}$  tăng\rq\rq\, sai.
  \itemch Xét dãy số $(u_n)$ có $u_n=\dfrac{1}{n}$. Với mọi $n\in \mathbb{N}^*$, ta có $\dfrac{u_{n+1}}{u_n} = \dfrac{n}{n+1}<1$, tức là $u_{n+1}<u_n$ cho nên dãy số $(u_n)$ đã cho  là dãy số giảm.\\
  Mệnh đề \lq\lq Dãy số $u_n=\dfrac{1}{n}$ không tăng\rq\rq\,  đúng.
  \itemch Xét  dãy số $(u_n)$ có $u_n=\dfrac{n+5}{3n+1}=\dfrac{1}{3}+\dfrac{14}{3(3n+1)}$.\\
  Với mọi $n\in \mathbb{N}^*$, ta có  $u_{n+1}-u_n = \dfrac{14}{3(3n+4)}-\dfrac{14}{3(3n+1)}<0$, tức là $u_{n+1}<u_n$ cho nên dãy số $(u_n)$ đã cho là dãy số giảm.\\
  Mệnh đề \lq\lq Dãy số $u_n=\dfrac{2n-1}{n+1}$  giảm\rq\rq\, đúng.
  \itemch Xét dãy số $(u_n)$  có $u_n=\dfrac{2n-1}{n+1}=2-\dfrac{3}{n+1}$.\\
  Với mọi $n\in \mathbb{N}^*$, ta có $u_{n+1}-u_n = \dfrac{3}{n+1}-\dfrac{3}{n+2}>0$, tức là $u_{n+1}>u_n$ cho  nên dãy số $(u_n)$ đã cho là dãy số tăng.\\
  Mệnh đề \lq\lq Dãy số $u_n=\dfrac{n+5}{3n+1}$ tăng\rq\rq\,  sai.
 \end{itemchoice}
 } 
\end{ex}

% nguồn: Ok_dungsai_lop11_Lan1
\begin{ex}%[1D2H1-4]%[ma:e398ca7f]
 Xét tính tăng giảm và bị chặn của dãy số $(u_n)$ với  $u_n=n^3+2n+1$. Các mệnh đề sau đúng hay sai?
  \choiceTF
   {Dãy số tăng, bị chặn}
  {Dãy số giảm, bị chặn  trên}
  {Dãy số giảm, bị  chặn}
  {\True Dãy số tăng,  bị chặn dưới}
  \loigiai{
   Ta có
  \[u_{n+1}-u_n=(n+1)^3+2(n+1)+1-(n^3+2n+1)=3n^2+3n+3>0,\,\,\forall n\in  \mathbb{N}^*.\]
  Nên  $(u_n)$ là dãy tăng.\\
  Mặt khác, vì $u_n>0,\,\,\forall n\in \mathbb{N^*}$  nên dãy $(u_n)$ bị chặn dưới.
  \begin{itemchoice}
   \itemch Mệnh  đề \lq\lq Dãy số tăng, bị chặn\rq\rq\, sai.
   \itemch Mệnh đề \lq\lq  Dãy số giảm, bị chặn trên\rq\rq\, sai.
   \itemch Mệnh đề \lq\lq Dãy số giảm, bị chặn\rq\rq\,  sai.
   \itemch  Mệnh đề \lq\lq Dãy số tăng, bị chặn dưới\rq\rq\, đúng.
  \end{itemchoice}
 } 
\end{ex}

% nguồn: Ok_dungsai_lop11_Lan1
\begin{ex}%[1D2H1-4]%[ma:10650e1b]
 Các mệnh đề sau đúng hay  sai?
  \choiceTF
  {Dãy số $(u_n)$ với $u_n=\dfrac{1}{2n^2+3}$, $\forall  n\in \mathbb{N}^*$ tăng}
  {\True Dãy số  $(u_n)$ với $u_n=(-1)^{2n}\cdot (n^2+1)$, $\forall n\in \mathbb{N}^*$ không tăng không giảm}
  {Dãy số $(u_n)$  với $u_n=\dfrac{1}{n(n+1)}$, $\forall n\in \mathbb{N}^*$ không giảm}
  {\True  Dãy số $(u_n)$ với $u_n=(-1)^n\cdot (n^2+4n+1)$, $\forall n\in \mathbb{N}^*$ giảm}
  \loigiai{
   Xét lần lượt các dãy số
  \begin{itemchoice}
   \itemch  $u_n=\dfrac{1}{2n^2+3}$\\
    Ta có
   \begin{eqnarray*}
    u_{n+1}-u_n&=&\dfrac{1}{2(n+1)^2+3}-\dfrac{1}{2n^2+3}=\dfrac{1}{2n^2+4n+5}-\dfrac{1}{2n^2+3}\\ 
     &=&\dfrac{-4n-2}{(2n^2+4n+5)(2n^2+3)}<0,\,\forall n\in\mathbb{N}^*.
   \end{eqnarray*}
   Do đó dãy số  này là dãy giảm.
   \itemch $u_n=(-1)^n\cdot (n^2+4n+1)$\\ 
   Dãy này  là dãy đan dấu nên không tăng không giảm.
    \itemch $u_n=\dfrac{1}{n(n+1)}$\\
   Ta  có $u_{n+1}-u_n=\dfrac{1}{(n+1)(n+2)}-\dfrac{1}{n(n+1)}=\dfrac{-2}{n(n+1)(n+2)}<0,\,\forall n\in \mathbb{N}^*$.\\
   Do đó  dãy số này là dãy giảm.
   \itemch $u_n=(-1)^{2n}\cdot (n^2+1)$\\ 
   Ta  có
   \begin{eqnarray*}
    u_{n+1}-u_n&=&(-1)^{2(n+1)}\cdot ((n+1)^2+1)-(-1)^{2n}\cdot (n^2+1)\\ 
    &=&n^2+2n+2-(n^2+1)=2n+1>0,\,\forall n\in \mathbb{N}^* 
   \end{eqnarray*}
   Do đó dãy  số này là dãy tăng.
  \end{itemchoice}
 } 
\end{ex}

% nguồn: Ok_dungsai_lop11_Lan1
\begin{ex}%[1D2H1-4]%[ma:7be239b3]
  Cho dãy số $ \left(u_n\right) $ với $ u_n=\dfrac{n+2}{2n+1}, \forall n\ge 1$. Các mệnh đề sau đúng hay sai?
  \choiceTF
  {\True Số $ \dfrac{5}{7} $  là số hạng thứ $3$ của dãy}
   {\True Dãy số $ u_n>0,\forall n\in\mathbb{N}^* $}
  {\True Dãy số $ \left(u_n\right)  $ là dãy số giảm}
   {Dãy số $ \left(u_n\right) $ là dãy số tăng}
  \loigiai{
  \begin{itemchoice}
   \itemch Ta có  $ \left(u_3\right) =\dfrac{5}{7}$ do đó mệnh đề \lq\lq Số $ \dfrac{5}{7} $ là số hạng thứ $3$ của dãy\rq\rq\, đúng.
   \itemch Mệnh  đề \lq\lq Dãy số $ u_n>0,\forall n\in\mathbb{N}^* $\rq\rq\, đúng.
    \itemch
   Ta có  $ u_{n+1}-u_n=\dfrac{n+3}{2n+3}-\dfrac{n+2}{2n+1}=\dfrac{-3}{(2n+1)(2n+3)}<0,\forall n\in\mathbb{N}^* $.\\
   Do đó mệnh  đề \lq\lq Dãy số $ \left(u_n\right) $ là dãy số giảm\rq\rq\, đúng.
   \itemch Mệnh đề \lq\lq Dãy số $ \left(u_n\right) $ là dãy số tăng\rq\rq\,  sai.
  \end{itemchoice}
  }
\end{ex}

% nguồn: Ok_dungsai_lop11_Lan1
\begin{ex}%[1D2H1-4]%[ma:f8aa20df]
 Các mệnh đề sau đúng hay  sai?
  \choiceTF
  {\True Dãy số $u_n = \dfrac {n + 3}{n + 1}$ giảm} 
  {Dãy số $u_n  = \cos n$ giảm}
  {\True Dãy số  $u_n = \dfrac{3^n}{n}$ tăng}
  {\True Dãy số $u_n = 25 - 10n - n^2$ không  tăng không giảm}
  \loigiai{
  \begin{itemchoice}
   \itemch Xét mệnh đề $u_n = \dfrac {n + 3}{n + 1}$, suy ra $u_n = 1 + \dfrac {2}{n + 1}$ khi $n$ tăng thì dãy số giảm , do đó mệnh đề \lq\lq Dãy số $u_n =  \dfrac {n + 3}{n + 1}$ giảm\rq\rq\ đúng.
   \itemch Dãy số $u_n = \cos n$ có $u_2 >u_3<u_4$ là dãy số không tăng không giảm nên mệnh đề \lq\lq Dãy số $u_n = \cos  n$ giảm\rq\rq\, sai.
   \itemch $\dfrac {u_{n + 1}}{u_n} = \dfrac{3^{n + 1}\cdot n}{3^n\cdot (n + 1)} = \dfrac {3n}{n + 1} = 3 - \dfrac{3}{n + 1} \ge 3 - \dfrac{3}{2} = \dfrac{3}{2} > 1$ ,$\forall n \ge 1$, do đó mệnh đề \lq\lq  Dãy số $u_n = \dfrac{3^n}{n}$ tăng\rq\rq\ đúng.
   \itemch Xét mệnh đề $u_n = 25 - 10n - n^2$, suy  ra $u_n = 25 - 10n - n^2 = - (n - 5)^2$ là dãy số không tăng không giảm, do đó mệnh đề \lq\lq Dãy số $u_n = 25 - 10n - n^2$ không tăng không giảm\rq\rq\ đúng.
  \end{itemchoice}
 } 
\end{ex}

% nguồn: Ok_dungsai_lop11_Lan2
\begin{ex}%[1D2H1-4]%[ma:b511bfe0]
	Cho dãy số $(u_n)$,  biết $u_n=1-\dfrac{1}{n}, \forall n \in \mathbb{N}^*$.
	 \choiceTF
		{\True Số hạng  $u_3=\dfrac{2}{3}$}
		{$u_7-u_8=\dfrac{1}{56}$} 
		{\True  $u_{n+1}-u_n=-\dfrac{1}{n(n+1) }$}
		{\True Dãy số $u_n$ là  dãy số tăng}
	 \loigiai{
		\begin{itemchoice}
			\itemch Ta có  $u_3=1-\dfrac{1}{3}=\dfrac{2}{3}$.
			 \itemch Số hạng $u_7-u_8=-\dfrac{1}{56}$.
			\itemch Ta có $u_{n+1}-u_n=\left(1-\dfrac{1}{n+1}\right) -\left(1-\dfrac{1}{n}\right) =\dfrac{1}{n}-\dfrac{1}{n+1}=\dfrac{1}{n(n+1) }>0$, $\forall  n \in \mathbb{N}^*$.\\
			Suy ra $u_{n+1}>u_n$, $\forall n \in \mathbb{N}^*$. \itemch Vậy dãy số $(u_n)$ là dãy  số tăng.
	\end{itemchoice}}
\end{ex}

% nguồn: Ok_dungsai_lop11_Lan2
\begin{ex}%[1D2H1-4]%[ma:2037cf81]
	Cho  dãy số $(u_n)$ có số hạng tổng quát $u_n=\sqrt{n+1}-\sqrt{n}$.
	 \choiceTF
		 {$\dfrac{u_{n+1}}{u_n}=\dfrac{\sqrt{n+2}+\sqrt{n}}{\sqrt{n+3}+\sqrt{n+2}}$}
		{\True $\dfrac{u_{2024}}{u_{2023}} < 1$} 
		{\True $u_{n+1}  < u_n$, $\forall n \in \mathbb{N}^*$}
		{\True Dãy số $(u_n)$ là dãy số  giảm}
	\loigiai{\begin{itemchoice}
		\itemch Nhận xét $u_n =  \sqrt{n+1} - \sqrt{n} > 0$, $\forall n \in \mathbb{N}^*$.\\
		Ta có\begin{eqnarray*}
		\dfrac{u_{n+1}}{u_n} &=& \dfrac{\sqrt{n+2} - \sqrt{n+1}}{\sqrt{n+1} -  \sqrt{n}}\\
		&=& \dfrac{\left(\sqrt{n+2} - \sqrt{n+1}\right) \left(\sqrt{n+2} + \sqrt{n+1}\right) \left(\sqrt{n+1} + \sqrt{n}\right) }{\left(\sqrt{n+1} - \sqrt{n}\right) \left(\sqrt{n+1} +  \sqrt{n}\right) \left(\sqrt{n+2} + \sqrt{n+1}\right) }\\
		&=& \dfrac{\sqrt{n+1} + \sqrt{n}}{\sqrt{n+2} +  \sqrt{n+1}}.
	\end{eqnarray*}
	\itemch Vì $0 < \sqrt{n+1} + \sqrt{n} < \sqrt{n+2} + \sqrt{n+1}$ nên $\dfrac{\sqrt{n+1} + \sqrt{n}}{\sqrt{n+2} + \sqrt{n+1}} < 1$  hay $\dfrac{u_{n+1}}{u_n} < 1$, $\forall n \in \mathbb{N}^*$.
	\itemch Suy ra  $u_{n+1} < u_n$, $\forall n \in \mathbb{N}^*$.
	\itemch Vậy dãy số $(u_n)$ là  dãy số giảm.
\end{itemchoice}}
\end{ex}

% nguồn: Ok_dungsai_lop11_Lan2
\begin{ex}%[1D2H1-4]%[ma:d52e3d4a]
	Cho dãy số $(u_n)$ có số  hạng tổng quát $u_n=1-\dfrac{1}{n}$.
	 \choiceTF
		{\True Ta có $u_3=\dfrac{2}{3}$} 
		{$u_7-u_8=\dfrac{1}{56}$} 
		{$u_{n+1}-u_n=-\dfrac{1}{n(n+1)  }$}
		{\True Dãy số $(u_n)$ là dãy số  tăng}
	\loigiai{\begin{itemchoice}
		\itemch Ta có  $u_3 = \dfrac{2}{3}$.
		\itemch $u_7-u_8  = -\dfrac{1}{56}$.
		\itemch Ta có $u_{n+1}-u_n = \left(1-\dfrac{1}{n+1}\right) -\left(1-\dfrac{1}{n}\right) = \dfrac{1}{n}-\dfrac{1}{n+1} = \dfrac{1}{n(n+1) }  > 0$, $\forall n \ge 1$.\\
		Suy ra $u_{n+1} > u_n$, $\forall n \in \mathbb{N}^*$. \itemch Vậy  dãy số $(u_n)$ là dãy số tăng.
\end{itemchoice}}
\end{ex}

% nguồn: Ok_dungsai_lop11_Lan2
\begin{ex}%[1D2H1-4]%[ma:d3aae649]
	Cho dãy số $\left(u_n \right)$ có số hạng  tổng quát $u_n=\dfrac{n}{4^n}$.
	\choiceTF 
		{Ta có $u_n=\dfrac{n}{4^n} < 0$, $\forall  n \in N^*$}
		{\True  Ta có $\dfrac{u_{n+1}}{u_n} < 1$, $\forall n \ge 1$}
		 {\True Ta có $u_{2024} < u_{2023}$}
		 {Dãy số $\left(u_n \right)$ là dãy số tăng}
	\loigiai{\begin{itemchoice}
		\itemch Ta có $u_n =  \dfrac{n}{4^n} > 0$, $\forall n \in N^*$.
		\itemch Ta có $\dfrac{u_{n+1}}{u_n} = \dfrac{n+1}{4^{n+1}}:\dfrac{n}{4^n} = \dfrac{n+1}{4n} = \dfrac{1}{4} + \dfrac{1}{4n} < 1$, $\forall  n \ge 1$.
		\itemch Suy ra $u_{n+1} < u_n$, $\forall n  \in N^*$.
		\itemch Vậy dãy số $\left(u_n \right)$  là dãy số giảm.
\end{itemchoice}}
\end{ex}

% nguồn: Ok_dungsai_lop11_Lan2
\begin{ex}%[1D2V1-4]%[ma:78071788]
	Cho dãy số $\left(u_n \right)$ biết $u_n=n^2+2n$,  $n \in \mathbb{N}^*$.
	\choiceTF 
		{\True Số  hạng đầu tiên của dãy số là $u_1=3$}
		 {Dãy số $\left(u_n \right)$ là một dãy số giảm}
		{Số $143$ là số  hạng thứ $13$ trong dãy số $\left(u_n \right)$}
		{Với mọi $n \in \mathbb{N}^*$ thì $\dfrac{1}{u_1}+\dfrac{1}{u_2}+\dfrac{1}{u_3}+\cdots+\dfrac{1}{u_n}=\dfrac{3n^2+5n}{2(n+1) (n+2)  }$}
	\loigiai{\begin{itemchoice}
		\itemch Ta có  $u_1 = 1^2 + 2\cdot 1 = 3$.
		\itemch $\forall n \in \mathbb{N}^*:u_{n+1} - u_n =(n+1)^2  + 2(n+1) -(n^2 + 2n) = 2n + 3 > 0$.\\
		Cho nên dãy số $\left(u_n \right)$ là một dãy số tăng. \itemch Ta có $u_n = 143 \Leftrightarrow n^2  + 2n = 143 \Leftrightarrow \hoac{&n=11 \\& n=-13.}$\\
		Do $n \in \mathbb{N}^*$  nên $n=11$, nghĩa là số $143$ là số hạng thứ $11$.
		 \itemch Ta có $u_n = n^2 + 2n = n(n+2)$.\\
		Do đó $\dfrac{1}{u_1} + \dfrac{1}{u_2} + \dfrac{1}{u_3} +\cdots+ \dfrac{1}{u_n} = \dfrac{1}{1\cdot3} + \dfrac{1}{2\cdot4} + \dfrac{1}{3\cdot 5} +\cdots+ \dfrac{1}{n(n+2)  }$.\\
		Phân tích$\dfrac{1}{1\cdot3} = \dfrac{1}{2} \left(\dfrac{1}{1} - \dfrac{1}{3} \right)$, $\dfrac{1}{2\cdot4} = \dfrac{1}{2} \left(\dfrac{1}{2} - \dfrac{1}{4}\right)$,$\cdots$, $\dfrac{1}{n(n+2) } = \dfrac{1}{2} \left(\dfrac{1}{n} - \dfrac{1}{n+2}  \right)$.\\
		Ta được $\dfrac{1}{u_1} + \dfrac{1}{u_2} + \dfrac{1}{u_3} +\cdots+ \dfrac{1}{u_n} = \dfrac{1}{2} \left(\dfrac{1}{1} - \dfrac{1}{2} +  \dfrac{1}{n+1} - \dfrac{1}{n+2} \right)$\\ $= \dfrac{1}{2}\cdot \dfrac{3(n+1) (n+2) - 2(n+2+n+1) }{2(n+1) (n+2) } = \dfrac{3n^2+5n}{4(n+1) (n+2) }$.
\end{itemchoice}}
\end{ex}


%%%%%%%%%%%%%%%% TRẢ LỜI NGẮN %%%%%%%%%%%%%%%%

% nguồn: Ok_TraLoiNgan_Lop11_Lan2
\begin{ex}%[1D2H1-4]%[ma:73eb20bb]
	Cho dãy số $(a_n)$ với $a_n=\dfrac{7n+5}{kn+7}$, ($k \in \mathbb{R}$). Giá trị nguyên lớn nhất của $k$  bằng bao nhiêu để $(a_n)$ là một dãy số tăng?
	\shortans{9} 
	\loigiai{Ta có $a_n = \dfrac{7n+5}{kn+7}$ và $a_{n+1} = \dfrac{7(n+1) +5}{k(n+1)  +7} = \dfrac{7n+12}{kn+k+7}$.\\
		Khi đó ta  có $a_{n+1} - a_n = \dfrac{7n+12}{kn+k+7} - \dfrac{7n+5}{kn+7} = \dfrac{-5k+49}{(kn+k+7) (kn+7) }$.\\
		Để dãy số tăng khi và chỉ khi $a_{n+1} - a_n >  0$, $\forall n\in \mathbb{N}^* \Leftrightarrow \dfrac{-5k+49}{(kn+k+7) (kn+7) } > 0 \Leftrightarrow k < \dfrac{49}{5}$.\\
	Vậy  số nguyên là lớn nhất là $k = 9$.}
\end{ex}

% nguồn: Ok_TraLoiNgan_Lop11_Lan1
\begin{ex}%[1D2V1-4]%[ma:76b60f04]
	Cho dãy số $\left(u_{n}\right)$ với $u_{n} =\dfrac{mn+1}{n+1}$. Có bao nhiêu giá trị nguyên của tham số $m$ thuộc $[0;2024]$ để dãy số $\left(u_{n} \right)$ là  dãy số tăng?
	\shortans{2023} 
	\loigiai{ 
		 Xét
		\begin{eqnarray*}
			&  u_{n+1} -u_{n} &=\dfrac{m(n+1)+1}{(n+1)+1} -\frac{mn+1}{n+1}\\
			&&=\dfrac{mn+m+1}{n+2} -\dfrac{mn+1}{n+1} \\ 
			&& =\dfrac{mn^{2} +2mn+m+n+1-\left(mn^{2} +2mn+n+2\right)}{(n+2)(n+1)} \\ 
			&& =\dfrac{m-1}{(n+2)(n+1)}. 
		\end{eqnarray*}
		Dãy số đã cho là dãy tăng $\Leftrightarrow \dfrac{m-1}{(n+2)(n+1)} >0,\forall n\in \mathbb{N}^{*}  \Leftrightarrow m>1$ \\(do $(n+2)(n+1)>0$, $\forall n \in \mathbb{N}^{*}$).\\
	 Vậy $m>1$ thỏa mãn đề bài. Do đó có $2023$ giá trị.}
\end{ex}

% nguồn: Ok_TraLoiNgan_Lop11_Lan1
\begin{ex}%[1D2V1-4]%[ma:4e9bc566]
	Cho dãy  số $(u_{n})$ với $\heva{&u_{1}=1; u_2=2\\&u_{n+2}=u_{n+1}+(1-a)u_{n}, \forall n \in \mathbb{N}^*}$.\\
	Có bao nhiêu giá trị nguyên của $a$ thuộc $\left[1;10\right]$ để dãy số $(u_{n})$ tăng? 
	 \shortans{9}
	 \loigiai{
		 Xét $u_{n+2}-u_{n+1}=u_{n+1}+(1-a)u_{n} -u_{n+1}=(a-1)(u_{n+1}-u_{n})$.\\
		Khi  đó
		 $$u_{3}-u_{2}=(a-1)(u_{2}-u_{1})=a-1$$
		$$u_{4}-u_{3}=(a-1)(u_{3}-u_{2})=(a-1)^2$$ 
		$$u_{5}-u_{4}=(a-1)(u_{4}-u_{3})=(a-1)^3.$$ 
		Suy ra  $u_{n+1}-u_{n}=(a-1)^{n-1}$.\\
		 Để dãy số $(u_{n})$ tăng thì $a-1>0 \Leftrightarrow a>1$.\\
	Vậy $a>1$ thì dãy số  $(u_{n})$ tăng.}
\end{ex}

% nguồn: Ok_TraLoiNgan_Lop11_Lan1
\begin{ex}%[1D2V1-4]%[ma:60f3a62b]
	Cho dãy số $(u_{n})$ với  $u_{n}=n-\sqrt{n^2-1}$. Có bao nhiêu giá trị nguyên của $m\in[3;9]$ để $u_m<1$?
	 \shortans{7}
	\loigiai{ 
		 Ta có $u_{n}=n-\sqrt{n^2-1}=\dfrac{(n-\sqrt{n^2-1})(n+\sqrt{n^2-1})}{n+\sqrt{n^2-1}}=\dfrac{1}{n+\sqrt{n^2-1}}$.\\
		Khi đó $u_{n+1}=\dfrac{1}{(n+1)+\sqrt{(n+1)^2-1}}  < \dfrac{1}{n+\sqrt{(n+1)^2-1}}$.\\
		Suy  ra $u_{n+1}<u_{n}$.\\
	Vậy dãy $u_{n}$ là  dãy giảm.}
\end{ex}

% nguồn: Ok_TraLoiNgan_Lop11_Lan1
\begin{ex}%[1D2V1-4]%[ma:dc5a18ff]
	Cho dãy số $(u_n)$ với $u_n=\dfrac{na+2}{n+1}$. Có bao nhiêu giá trị của $a\in[-2024;2025]$  để dãy $(u_n)$ là dãy số tăng?
	 \shortans{2023}
	 \loigiai{
		Ta có $u_n=\dfrac{n a+2}{n+1}=a+\dfrac{2-a}{n+1} \Rightarrow u_{n+1}=a+\dfrac{2-a}{n+2} \Rightarrow  u_{n+1}-u_n=\dfrac{a-2}{(n+1)(n+2)}$.\\
		Để $\left(u_n\right)$ là dãy số tăng thì  $u_{n+1}-u_n=\dfrac{a-2}{(n+1)(n+2)}>0 \Leftrightarrow a>2,\forall n \in \mathbb{N}^*$.\\
		Vì $a\in[-2024;2025]$ nên $a\in\left\{3;4;\ldots 2025\right\}$.\\ 
		Vậy có $2023$  giá trị của $a$ thỏa mãn.
	} 
\end{ex}

% nguồn: Ok_TraLoiNgan_Lop11_Lan2
\begin{ex}%[1D2V1-4]%[ma:fc30ef4c]
	Cho dãy số $(u_n)$ biết $\heva{&u_1=1,  u_2=2\\&u_{n+2}=au_{n+1}+(1-a)u_n}$, $\forall n\in\mathbb{N}^*$. Số các giá trị nguyên của $a$ trong khoảng $(-2\,026;2\,026)$ để $(u_n)$ là dãy tăng
	\shortans{2,02} 
	 \loigiai{
		Xét hiệu 
		\begin{eqnarray*}
			 u_{n+2}-u_{n+1}&=&au_{n+1}+(1-a)u_n-u_{n+1}=(a-1)(u_{n+1}-u_n)\\
			\Rightarrow u_3-u_2&=&(a-1)(u_2-u_1)=a-1\\ 
			 \Rightarrow u_4-u_3&=&(a-1)(u_3-u_2)=(a-1)^2\\
			&\ldots&\\ 
			 \Rightarrow u_{n+1}-u_n&=&(a-1)^{n-1}>0.
		\end{eqnarray*}
		Để dãy số $(u_n)$ là dãy tăng thì  $a>1$.
	 }
\end{ex}

% nguồn: Ok_TraLoiNgan_Lop11_Lan2
\begin{ex}%[1D2V1-4]%[ma:ff2ebad5]
	Cho dãy số $(u_n)$ với $u_n=\dfrac{an^2+2}{n^2+3}$. Có bao nhiêu số nguyên $a\in[-10;10]$ để $(u_n)$ là dãy  số giảm?
	 \shortans{11}
	 \loigiai{
		Ta  có
		$u_n=\dfrac{an^2+2}{n^2+3}=a+\dfrac{2-3a}{n^2+3},n\in\mathbb{N}^*$.\\ 
		Xét 
		$u_{n+1}-u_n=\dfrac{2-3a}{(n+1)^2+3}-\dfrac{2-3a}{n^2+3}=(2-3a)\left[\dfrac{1}{(  n+1)^2+3}-\dfrac{1}{n^2+3} \right]$.\\
		Ta có $\dfrac{1}{(n+1)^2+3}<\dfrac{1}{n^2+3}$ nên $(u_n)$ là dãy số giảm khi và  chỉ khi $2-3a>0\Leftrightarrow a<\dfrac{2}{3}$.\\
		Các số nguyên $a\in \left[-10;10 \right]$ thỏa mãn là $ a\in  \lbrace -10; -9; -8; \ldots; 0 \rbrace$.\\
		Vậy có  $11$ số.
	 }
\end{ex}

% nguồn: Ok_TraLoiNgan_Lop11_Lan2
\begin{ex}%[1D2V1-4]%[ma:93aabcbc]
	Cho dãy số $(u_n)$ với $u_n=\dfrac{na+3}{2n+3}$. Số các giá trị nguyên của $a$ trong  khoảng $(-2\,023;2\,023)$ để $(u_n)$ là dãy số giảm là
	 \shortans{2,02}
	\loigiai{ 
		 Ta có $u_n=\dfrac{na+3}{2n+3}=\dfrac{a}{2}+\dfrac{3-\dfrac{3a}{2}}{2n+3}$.\\
		$u_{n+1}=\dfrac{a}{2}+\dfrac{3-\dfrac{3a}{2}}{2(n+1)+3}=\dfrac{a}{2}+\dfrac{3-\dfrac{3a}{2}}{2n+5}$.\\ 
		Xét 
		\begin{eqnarray*}
			 u_{n+1}-u_n=\dfrac{3-\dfrac{3a}{2}}{2n+5}-\dfrac{3-\dfrac{3a}{2}}{2n+3}&=&\left(3-\dfrac{3a}{2}\right)\left(\dfrac{1}{2n+5}-\dfrac{1}{2n+3}\right)\\&=&\left(3-\dfrac{3a}{2}\right)\cdot\dfrac{-2}{(2n+5)(2n+3)}
		\end{eqnarray*}
		Để $(u_n)$ là dãy  số giảm thì $u_{n+1}-u_n<0\Leftrightarrow 3-\dfrac{3a}{2}>0\Leftrightarrow\dfrac{3a}{2}<3\Leftrightarrow a<2$.\\
		Mà $a$ nguyên và $a$  thuộc $(-2\,023;-2\,023)$ nên $a\in\{-2\,022;-2\,021;\ldots;0;1\}$.\\
		Vậy có $1-(-2\,022)+1=2\,024$ giá trị  nguyên của $a$.
	 }
\end{ex}

% nguồn: Ok_TraLoiNgan_Lop11_Lan2
\begin{ex}%[1D2V1-4]%[ma:413f9d15]
	Cho dãy số $(u_n)$ biết $u_n=\dfrac{2an^2+1}{n^2+3}$, $\forall n\in \mathbb{N}^*$ và $a\in \mathbb{R}$. Có bao nhiêu giá trị nguyên của $a\in[-10;10]$ sao cho $(u_n)$ là một dãy số  giảm?
	 \shortans{11}
	\loigiai{ 
		Ta có  $u_n=2a+\dfrac{1-6a}{n^2+3}$.\\
		Xét hiệu $u_{n+1}-u_n=2a+\dfrac{1-6a}{(n+1)^2+3}-2a+\dfrac{1-6a}{n^2+3}=\left( 1-6a \right)\left[\dfrac{1}{{{\left( n+1  \right)}^2}+3}-\dfrac{1}{n^2+3} \right]$.\\
		Vì 
		\begin{eqnarray*}
			(n+1)^2  &>& n^2\\
			(n+1)^2 + 3 &>& n^2  + 3\\
			\dfrac{1}{(n+1)^2 + 3} &<& \dfrac{1}{n^2 +  3}\\
			\dfrac{1}{(n+1)^2 + 3} &-& \dfrac{1}{n^2 + 3} <  0
		\end{eqnarray*}
		$(u_n)$ là dãy số giảm khi và chỉ khi $u_{n+1}-u_n<0,\forall  n\in {{\mathbb{N}}^{*}}$
		 $\Leftrightarrow 1-6a>0\Leftrightarrow a<\dfrac{1}{6}$.\\
		Mà $a\in [-10;10]$ và $a$ là  số nguyên nên $a\in \left\{ -10;-9;-8;\ldots;0 \right\}$.
		Vậy có $11$ giá trị $a$. 
	} 
\end{ex}

% nguồn: Ok_TraLoiNgan_Lop11_Lan2
\begin{ex}%[1D2V1-4]%[ma:72ea70f1]
	Gọi $A$ là tập hợp tất cả giá trị nguyên của $a$ thuộc đoạn $[-5;5]$ sao cho dãy số $(u_n)$ với $u_n=\dfrac{an^4+2}{2n^4+5}$, $\forall n\in\mathbb{N}^*$ là một dãy số tăng. Hỏi tập  hợp $A$ có bao nhiêu phần tử?
	\shortans{5} 
	\loigiai{ 
		 Ta có $u_n=\dfrac{an^4+2}{2n^4+5}=\dfrac{a}{2}+\dfrac{4-5a}{2(2n^4+5)}$, $\forall n\in\mathbb{N}^*$.\\
		Xét 
		\begin{eqnarray*}
			u_{n+1}-u_n&=&\dfrac{4-5a}{2\left[2(n+1)^4+5\right]}-\dfrac{4-5a}{2(2n^4+5)}\\&=&\dfrac{4-5a}{2}\left[\dfrac{1}{2(n+1)^4+5}-\dfrac{1}{2n^4+5}\right]\\ 
			&=&\dfrac{4-5a}{2}\cdot\dfrac{2n^4+5-2(n+1)^4-5}{\left[2(n+1)^4+5\right]\left(2n^4+5\right)}\\ 
			 &=&(4-5a)\cdot\dfrac{n^4-(n+1)^4}{\left[2(n+1)^4+5\right]\left(2n^4+5\right)}
		\end{eqnarray*}
		Mà $\dfrac{n^4-(n+1)^4}{\left[2(n+1)^4+5\right]\left(2n^4+5\right)}<0$,  $\forall n\in\mathbb{N}^*$.\\
		Nên $(u_n)$  tăng $\Leftrightarrow u_{n+1}-u_n>0,\forall n\in\mathbb{N}^*\Leftrightarrow 4-5a<0\Leftrightarrow a>\dfrac{4}{5}$.\\
		Mà $a \in [-5;5] $ nên $A  = \{1;2;3;4;5\}$. Vậy $A$ có 5 phần tử.
	} 
\end{ex}

% nguồn: Ok_TraLoiNgan_Lop11_Lan2
\begin{ex}%[1D2V1-4]%[ma:04be742a]
	Cho dãy số $(u_n)$ biết $u_n=(2-a)n+a-2$, $\forall n\in\mathbb{N}^*$. Có bao nhiêu giá trị nguyên dương của $a$  để dãy số $u_n$ là dãy tăng?
	\shortans{1} 
	 \loigiai{
		Ta có $u_{n+1}=(2-a)(n+1)+a-2$.\\ 
		Suy ra $u_{n+1}-u_n=(2-a)(n+1)+a-2-\left[(2-a)n+a-2\right]=2-a$,  $\forall n\in\mathbb{N}^*$.\\
		$(u_n)$ là dãy tăng $\Leftrightarrow u_{n+1}-u_n>0$, $\forall n\in\mathbb{N}^*\Leftrightarrow 2-a>0\Leftrightarrow a<2$. 
	 }
\end{ex}

% nguồn: Ok_TraLoiNgan_Lop11_Lan2
\begin{ex}%[1D2V1-4]%[ma:0cbc20d3]
	Cho dãy số $\left(u_n \right)$ biết $\heva{&u_1=1;u_2=2 \\& u_{n+2}=au_{n+1}+(1-a) u_n, \forall n \in \mathbb{N}^*.}$ Giá trị  nguyên dương nhỏ nhất của $a$ bằng bao nhiêu để $(u_n)$ là dãy số tăng?
	 \shortans{2}
	\loigiai{Xét hiệu $u_{n+2} -  u_{n+1} = au_{n+1} +(1-a) u_n - u_{n+1} =(a-1) u_{n+1} +(1-a) u_n =(a-1) \left(u_{n+1} - u_n \right)$.\\
		Khi đó: $u_3 - u_2 =(a-1) \left(u_2 - u_1  \right) =(a-1) (2-1) = a-1$.\\
		$u_4 - u_3 =(a-1) \left(u_3 -  u_2 \right) =(a-1) (a-1) = (a-1)^2$.\\
		$u_5 - u_4 =(a-1) \left(u_4 - u_3 \right) =(a-1) (a-1)^2 = (a-1)^3$.\\ 
		Suy ra $u_{n+1} - u_n = (a-1)^{n-1}$.\\ 
		Để dãy số  $\left(u_n \right)$ tăng thì $a-1 > 0 \Leftrightarrow a > 1$.\\
		 Suy ra $a > 1$ thì dãy số $\left(u_n \right)$ tăng.\\
	Vậy  số nguyên dương $a$ nhỏ nhất để $(u_n)$ là dãy tăng là $a=2$. }
\end{ex}

% nguồn: Ok_TraLoiNgan_Lop11_Lan2
\begin{ex}%[1D2V1-4]%[ma:5944daef]
	Cho dãy số $(u_n)$  biết $\heva{&u_1=3\\&u_{n+1}=2au_n+3,\forall n\geq1}$. Có bao nhiêu giá trị nguyên của $a$ trong khoảng $(-2\,025;2\,025)$ để $(u_n)$ là dãy số tăng?
	 \shortans{2024}
	\loigiai{ 
		Ta  có
		\begin{eqnarray*}
			 u_{n+1}-u_n &=&(2au_n+3)+(2au_{n-1}+3)=2a\left(u_n-u_{n-1}\right),\forall n\geq2 \\
			\Rightarrow u_3-u_2&=&2a(u_2-u_1)  = 2a(6a+3-3)=2a\cdot6a=3(2a)^2\\
			\Rightarrow u_4-u_3&=&2a(u_3-u_2)=2a\cdot3\cdot(2a)^2=3\cdot(2a)^2\\ 
			\ldots&&\\ 
			\Rightarrow u_{n+1}-u_n &=&3(2a)^n.\\ 
		\end{eqnarray*}
		Để dãy số $(u_n)$ là dãy số  tăng thì $u_{n+1}-u_n=3(2a)^n>0$, $\forall n\geq1$.\\
		Suy ra  $a>0$.
	} 
\end{ex}


%%%%%%%%%%%%%%%% TỰ LUẬN %%%%%%%%%%%%%%%%

% nguồn: Ok_TuanLuan_Toan11_Lan2
\begin{ex}%[1D2N1-4]%[ma:836fb454]
	Xét tính đơn điệu  của dãy số $(u_n)$ biết $u_n=3n+6$.
	 \loigiai{
		Ta có $u_n=3n+6\Rightarrow  u_{n+1}=3(n+1)+6=3n+9$.\\
		Xét hiệu $u_{n+1}-u_n=(3n+9)-(3n+6)=3> 0,\  \forall n \in \mathbb{N}^*$.\\
		Vậy  $(u_n)$ là dãy số tăng.\\
		\textbf{Giải nhanh:} Dãy này có dạng  $u_n=a n+b; a=3> 0$ nên dãy số tăng
	} 
\end{ex}

% nguồn: Ok_TuanLuan_Toan11_Lan2
\begin{ex}%[1D2H1-4]%[ma:a0aaf3c4]
	Xét  tính đơn điệu của dãy số $(u_n)$ biết $u_n=\dfrac{n+5}{n+2}$.
	 \loigiai{
		Ta có $u_n=\dfrac{n+5}{n+2}=1+\dfrac{3}{n+2} \Rightarrow  u_{n+1}=1+\dfrac{3}{n+3}$.\\
		Xét  hiệu $u_{n+1}-u_n=\dfrac{3}{n+3}-\dfrac{3}{n+2}=\dfrac{-3}{(n+2)(n+3)} < 0,\ \forall n \in \mathbb{N}^*$.\\
		Vậy $(u_n)$ là dãy số  giảm.\\
		\textbf{Giải nhanh:} Dãy này  có dạng $u_n=\dfrac{an+b}{cn+d}$.\\
		Mẫu $n+2>  0,\ \forall n \in \mathbb{N}^*$ và $ad-bc=2-5=-3<0$ nên $(u_n)$ là dãy số giảm.
	 }
\end{ex}

% nguồn: Ok_TuanLuan_Toan11_Lan2
\begin{ex}%[1D2H1-4]%[ma:70730a4f]
	Xét tính  tăng, giảm của dãy số $(u_n)$ biết $u_n=\dfrac{5^n}{n^2}$.
	\loigiai{ 
		 Ta có $u_n=\dfrac{5^n}{n^2} > 0, \forall n \in \mathbb{N}^* \Rightarrow u_{n+1}=\dfrac{5^{n+1}}{(n+1)^2}$.\\
		Xét tỉ số 
		\begin{eqnarray*}
			\dfrac{u_{n+1}}{u_n}&=&\dfrac{5^{n+1}}{(n+1)^2} \cdot  \dfrac{n^2}{5^n}\\
			 &=&\dfrac{5n^2}{n^2+2n+1}\\
			 &=&\dfrac{n^2+2n+1+4n^2-2n-1}{n^2+2n+1}\\
			&=&1+\dfrac{2n(n-1)+2n^2-1}{n^2+2 n+1}>1,\ \forall  n \in \mathbb{N}^*.
		\end{eqnarray*}
		Vậy  $(u_n)$ là dãy số tăng.
	} 
\end{ex}

% nguồn: Ok_TuanLuan_Toan11_Lan2
\begin{ex}%[1D2H1-4]%[ma:0c0bcb13]
	Với giá trị nào của $a$ thì  dãy số $(u_n)$, với $u_n=\dfrac{n a+2}{n+1}$
	\begin{enumerate}
		\item  là dãy số tăng.
		\item là  dãy số giảm.
	\end{enumerate}
	\loigiai{ 
		Ta có $u_n=\dfrac{n a+2}{n+1}=a+\dfrac{2-a}{n+1} \Rightarrow u_{n+1}=2+\dfrac{2-a}{n+2} \Rightarrow  u_{n+1}-u_n=\dfrac{a-2}{(n+1)(n+2)}$.
		\begin{enumerate}
			\item Để $(u_n)$ là dãy số tăng thì $u_{n+1}-u_n=\dfrac{a-2}{(n+1)(n+2)} > 0\Leftrightarrow a  > 2$.
			\item Để $(u_n)$ là dãy số giảm thì $u_{n+1}-u_n=\dfrac{a-2}{(n+1)(n+2)} < 0\Leftrightarrow a < 2$. 
	\end{enumerate}}
\end{ex}

% nguồn: Ok_TuanLuan_Toan11_Lan2
\begin{ex}%[1D2H1-4]%[ma:2d9ecf21]
	Xét tính đơn điệu của dãy  số $(u_n)$ biết
	\begin{enumerate}
		\item $u_n=\dfrac{5^n}{n^2}$; 
		\item $\heva{&u_1=2\\&u_n=\dfrac{3u_{n-1}+1}{4}, \forall n  \geq 2.}$
	\end{enumerate}
	\loigiai{\begin{enumerate}
		\item Ta có $u_n=\dfrac{5^n}{n^2} > 0,  \forall n \in \mathbb{N}^* \Rightarrow u_{n+1}=\dfrac{5^{n+1}}{(n+1)^2}$.
		 \\
		Xét tỉ  số
		 \allowdisplaybreaks
		\begin{eqnarray*}
			\dfrac{u_{n+1}}{u_n}&=&\dfrac{5^{n+1}}{(n+1)^2} \cdot \dfrac{n^2}{5^n}\\ 
			 &=&\dfrac{5n^2}{n^2+2n+1}\\
			 &=&\dfrac{n^2+2n+1+4n^2-2n-1}{n^2+2n+1}\\
			&=&1+\dfrac{2n(n-1)+2n^2-1}{n^2+2n+1} >  1, \forall n \in \mathbb{N}^*.
		\end{eqnarray*}
		Vậy $(u_n)$ là dãy  số tăng.
		\item Ta dự đoán  dãy số giảm sau đó ta sẽ chứng minh nó giảm.\\
		 Ta có $u_n-u_{n-1}=\dfrac{3u_{n-1}+1}{4}-u_{n-1}=\dfrac{1-u_{n-1}}{4}$.\\
		Do đó, để chứng minh dãy $(u_n)$ giảm ta chứng minh $u_n > 1$,  $\forall n \geq 1$ bằng phương pháp quy nạp toán học.\\
		Thật vậy, với $n=1\Rightarrow  u_1=2> 1$.\\
		Giả sừ $u_k  > 1\Rightarrow u_{k+1}=\dfrac{3u_k+1}{4} > \dfrac{3+1}{4}=1$.\\
		Theo nguyên lí quy  nạp ta có $u_n > 1\forall n \geq 1$.\\
		Suy ra $u_n-u_{n-1} < 0\Leftrightarrow u_n <  u_{n-1} \forall n \geq 2$ hay dãy $(u_n)$ giảm.
	\end{enumerate}
} 
\end{ex}

% nguồn: Ok_TuanLuan_Toan11_Lan2
\begin{ex}%[1D2V1-4]%[ma:efca5f87]
	Cho dãy  số $(u_n)$ biết $(u_n) \colon \heva{&u_1=2\\ &u_n=\dfrac{3u_{n-1}+1}{4} \forall n \geq 2}$. Xét tính tăng, giảm của dãy số $(u_n)$.
	 \loigiai{
		Ta dự đoán dãy số giảm sau đó ta sẽ chứng minh  nó giảm.\\
		 Ta có $u_n-u_{n-1}=\dfrac{3u_{n-1}+1}{4}-u_{n-1}=\dfrac{1-u_{n-1}}{4}$.\\
		Do đó, để chứng  minh dãy $(u_n)$ giảm ta chứng minh $u_n > 1,\ \forall n \geq 1$ bằng phương pháp quy nạp toán học. Thật vậy\\
		Với $n=1\Rightarrow u_1=2>  1$.\\
		Giả sử $u_k > 1\Rightarrow  u_{k+1}=\dfrac{3u_k+1}{4} > \dfrac{3+1}{4}=1$.\\
		Theo nguyên lí quy nạp  ta có $u_n > 1,\ \forall n \geq 1$.\\
		Suy ra $u_n-u_{n-1} < 0\Leftrightarrow u_n < u_{n-1},\ \forall n  \geq 2$ hay dãy ($u_n$) giảm.\\
		\textbf{Giải nhanh:} Dãy $(u_n)$  có dạng $u_{n+1}=a u_n+b$.\\
		Ở đây $a=\dfrac{3}{4} >  0$ và $u_2-u_1=\dfrac{7}{4}-2=-\dfrac{1}{4} < 0$.\\
		Suy ra dãy số  giảm.\\
		Tổng quát ta có thể chứng  minh dãy số $(u_n)\colon \heva{&u_1=c>1\\ &u_n=\dfrac{au_{n-1}+b}{a+b},(a,b>0),\ \forall n \geq 2}$ giảm tương tự như trên.
	 }
\end{ex}

% nguồn: Ok_TuanLuan_Toan11_Lan2
\begin{ex}%[1D2V1-4]%[ma:254a8807]
	Xét tính tăng giảm của  các dãy số sau, biết
	\begin{multicols}{3}
		\begin{enumerate}
			\item $u_n=2n+3$. 
			 \item $u_n=\dfrac{\sqrt{n}}{2^n}$.
			\item  $u_n=\dfrac{n}{n^2+1}$.
			\item  $u_n=\dfrac{\sqrt{n+1}-n}{n}$.
			\item $u_n=\dfrac{1}{n}-2$. 
			\item $u_n=\dfrac{n-1}{n+1}$. 
			 \item $u_n=\dfrac{2n+1}{5n+2}$.
			 \item $u_n=2n^2+5$.
			\item  $u_n=\dfrac{2n^2-1}{n^2+1}$.
			\item $u_n=\sqrt{n+1}-\sqrt{n}$. 
			\item $u_n=\dfrac{3n^2-2n+1}{n+1}$. 
			\item $u_n=\dfrac{\sqrt{n+1}-1}{n}$. 
		\end{enumerate}
	\end{multicols}
	 \loigiai{
		\begin{enumerate}
			\item Ta có $u_n=2n+3; u_{n+1}=2(n+1)+3=2n+5\Rightarrow u_{n+1}-u_n=(2n+5)-(2n+3) > 0$,  $\forall n\in\mathbb{N}^*$.\\
			Suy ra  $u_{n+1} > u_n$, $\forall n\in\mathbb{N}^*$.\\
			Vậy  dãy số đã cho là dãy tăng.
			\item Ta có  $u_n=\dfrac{\sqrt{n}}{2^n}; u_{n+1}=\dfrac{\sqrt{n+1}}{2^{n+1}} \Rightarrow \dfrac{u_{n+1}}{u_n}=\dfrac{\sqrt{n+1}}{2^{n+1}} \cdot \dfrac{2^n}{\sqrt{n}}=\frac{1}{2} \dfrac{\sqrt{n+1}}{\sqrt{n}}=\frac{1}{2} \sqrt{\dfrac{n+1}{n}}$.\\
			Giả sừ $\dfrac{u_{n+1}}{u_n}=\frac{1}{2} \sqrt{\dfrac{n+1}{n}} >  1\Leftrightarrow \frac{1}{4} \dfrac{n+1}{n} > 1\Leftrightarrow n+1> 4n \Leftrightarrow 3n < 1\Rightarrow$ vô lý.\\
			Vậy $\dfrac{u_{n+1}}{u_n} <  1\Leftrightarrow u_{n+1} < u_n$, $\forall n\in\mathbb{N}^*$.\\
			Vậy dãy số đã cho là  dãy số giảm.
			 \item Ta có $u_n=\dfrac{n}{n^2+1}$; $u_{n+1}=\dfrac{n+1}{(n+1)^2+1}=\dfrac{n+1}{n^2+2n+2}$.
			 Khi đó
			 \allowdisplaybreaks
			\begin{eqnarray*}
				u_{n+1}-u_n&=&\dfrac{n+1}{n^2+2n+2}-\dfrac{n}{n^2+1}\\ 
				&=&\dfrac{(n+1)\left(n^2+1\right)-n\left(n^2+2n+2\right)}{\left(n^2+1\right)\left(n^2+2n+2\right)}\\ 
				&=&\dfrac{n^3+n^2+n+1-n^3-2n^2-2n}{\left(n^2+1\right)\left(n^2+2n+2\right)}\\ 
				&=&\dfrac{-n^2-n+1}{\left(n^2+1\right)\left(n^2+2n+2\right)} < 0,\forall n  \geq 1.
			\end{eqnarray*}
			Vậy $(u_n)$ là dãy  số giảm.
			\item Ta có $u_n=\dfrac{\sqrt{n+1}-n}{n}=\dfrac{\sqrt{n+1}}{n}-1 \Rightarrow u_{n+1}=\dfrac{\sqrt{n+2}}{n+1}-1$.\\ 
			Khi đó \allowdisplaybreaks 
			\begin{eqnarray*}
				u_{n+1}-u_n&=&\left(\dfrac{\sqrt{n+2}}{n+1}-1\right)-\left(\dfrac{\sqrt{n+1}}{n}-1\right)\\ 
				 &=&\dfrac{\sqrt{n+2}}{n+1}-\dfrac{\sqrt{n+1}}{n}\\
				&=&\dfrac{n \sqrt{n+2}-(n+1)  \sqrt{n+1}}{n(n+1)}.
			\end{eqnarray*}
			 Giả sử
			\allowdisplaybreaks 
			\begin{eqnarray*}
				u_{n+1}-u_n>0 &\Leftrightarrow&  n \sqrt{n+2}-(n+1) \sqrt{n+1}>0\\
				& \Leftrightarrow&  n \sqrt{n+2}>(n+1) \sqrt{n+1}\\
				&\Leftrightarrow&  n^2(n+2)>(n+1)^3\\
				& \Leftrightarrow& n^3+2 n^2>n^3+3 n^2+3 n+1\\ &\Leftrightarrow &n^2+3  n+1<0 ~(\text { vô lý}).
			\end{eqnarray*}
			Suy ra $u_{n+1}-u_n<0$, $\forall n\in\mathbb{N}^*$,  vậy $\left(u_n\right)$ là dãy số giảm.
			 \item $u_n=\dfrac{1}{n}-2 \Rightarrow u_{n+1}=\dfrac{1}{n+1}-2$.\\
			 Suy ra $u_{n+1}-u_n=\left(\dfrac{1}{n+1}-2\right)-\left(\dfrac{1}{n}-2\right)=-\dfrac{1}{n(n+1)}<0$, $\forall n\in\mathbb{N}^*$.\\
			Suy ra $u_{n+1}<u_n$, $\forall  n\in\mathbb{N}^*$. Vậy dãy số $\left(u_n\right)$ là dãy số giảm.
			\item Ta có $u_n=\dfrac{n-1}{n+1}=1-\dfrac{2}{n+1}$, suy ra $u_{n+1}=1-\dfrac{2}{n+2}$.\\ 
			 \allowdisplaybreaks
			\begin{eqnarray*}
				 u_{n+1}-u_n&=&\left(1-\dfrac{2}{n+2}\right)-\left(1-\dfrac{2}{n+1}\right)\\
				 &=&\dfrac{2}{(n+1)(n+2)}>0.
			\end{eqnarray*}
			Suy ra $u_{n+1}>u_n$, $\forall n\in\mathbb{N}^*$.\\ 
			Vậy dãy số $\left(u_n\right)$ là dãy  số tăng.
			\item Ta có $u_n=\dfrac{2  n+1}{5 n+2}=\dfrac{2}{3}+\dfrac{1}{5(5 n+2)} \Rightarrow u_{n+1}=\dfrac{2}{5}+\dfrac{1}{5(5 n+7)}$.\\
			Khi đó $u_{n+1}-u_n=\left(\dfrac{2}{5}+\dfrac{1}{5(5 n+7)}\right)-\left(\dfrac{2}{5}+\dfrac{1}{5(5 n+2)}\right)=-\dfrac{1}{(5 n+2)(5 n+7)}<0$, $\forall n\in\mathbb{N}^*$.\\ 
			Suy ra $u_{n+1}<u_n$, $\forall  n\in\mathbb{N}^*$.\\
			Vậy $\left(u_n\right)$ là dãy số  giảm.
			\item  $u_n=2 n^2+5 \Rightarrow u_{n+1}=2(n+1)^2+5$.\\
			Khi đó  $u_{n+1}-u_n=2(n+1)^2+5-\left(2 n^2+5\right)=4 n+2>0$, $\forall n\in\mathbb{N}^*$.\\
			$\Rightarrow u_{n+1}>u_n$, $\forall n\in\mathbb{N}^*$.\\ Vậy $\left(u_n\right)$ là dãy  số tăng.
			 \item Ta có $u_n=\dfrac{2 n^2-1}{n^2+1}=2-\dfrac{3}{n^2+1} \Rightarrow u_{n+1}=2-\dfrac{3}{(n+1)^2+1}$.\\
			Với $n \in \mathbb{N}^*$,  ta có
			 \allowdisplaybreaks
			\begin{eqnarray*}
				(n+1)^2>n^2 &\Leftrightarrow&  \dfrac{1}{(n+1)^2+1}<\dfrac{1}{n^2+1} \\
				&\Leftrightarrow& 2-\dfrac{3}{(n+1)^2+1}>2-\dfrac{3}{n^2+1}\\ &\Leftrightarrow& u_{n+1}>u_n. 
			\end{eqnarray*}
			$\Rightarrow\left(u_n\right)$ là dãy số tăng. 
			\item $u_n=\sqrt{n+1}-\sqrt{n}=\dfrac{1}{\sqrt{n}+\sqrt{n+1}} \Rightarrow u_{n+1}=\dfrac{1}{\sqrt{n+1}+\sqrt{n+2}}$.\\ 
			Với $n \in \mathbb{N}^*$, ta có 
			\allowdisplaybreaks 
			\begin{eqnarray*}
				&&\sqrt{n+2}+\sqrt{n+1}>\sqrt{n+1}+\sqrt{n}\\ 
				&\Rightarrow&  u_{n+1}=\dfrac{1}{\sqrt{n+2}+\sqrt{n+1}}<u_n=\dfrac{1}{\sqrt{n+1}+\sqrt{n}}.
			\end{eqnarray*}
			 $\Rightarrow u_{n+1}<u_n \Rightarrow\left(u_n\right)$ là dãy số giảm.
			\item $u_n=\dfrac{3 n^2-2 n+1}{n+1}=3 n-5+\dfrac{6}{n+1} \Rightarrow u_{n+1}=3 n-2+\dfrac{6}{n+1}$.\\ 
			Khi  đó $u_{n+1}-u_n=3 n-2+\dfrac{6}{n+2}-\left(3 n-5+\dfrac{6}{n+1}\right)=3-\dfrac{6}{(n+1)(n+2)}$.\\
			Với $n  \in \mathbb{N}^*$, ta có
			 \allowdisplaybreaks
			\begin{eqnarray*}
				(n+1)(n+2) \geq 6 &\Leftrightarrow& \dfrac{6}{(n+1)(n+2)} \leq 1  \\
				&\Leftrightarrow& 3-\dfrac{6}{(n+1)(n+2)} \geq  2\\
				 &\Rightarrow& u_{n+1}>u_n.
			\end{eqnarray*}
			$\Rightarrow\left(u_n\right)$  là dãy số tăng.
			\item Ta có $u_n=\dfrac{\sqrt{n+1}-1}{n}=\dfrac{n}{n(\sqrt{n+1}+1)}=\dfrac{1}{\sqrt{n+1}+1}$.\\ 
			Khi $n$ tăng thì dễ  thấy mẫu số tăng, phân số giảm nên dãy số đã cho là dãy số giảm.
		\end{enumerate}
	} 
\end{ex}
